\documentclass[10pt,letterpaper]{article}
\usepackage[english]{babel}
\usepackage[T1]{fontenc}
\usepackage[utf8]{inputenc}
\usepackage{lmodern,microtype}
\usepackage[margin=0.73in,headheight=15pt]{geometry}
\usepackage{graphicx,xcolor,booktabs,longtable,array}
\usepackage{listings,tcolorbox,titlesec,fancyhdr,hyperref,needspace}
\definecolor{ink}{HTML}{17253C}
\definecolor{accent}{HTML}{C85E32}
\definecolor{muted}{HTML}{536276}
\definecolor{wash}{HTML}{F1F5F8}
\definecolor{codebg}{HTML}{F7F9FB}
\hypersetup{colorlinks=true,linkcolor=ink,urlcolor=accent,pdftitle={ICPC Notebook Java Team Reference}}
\titleformat{\section}{\Large\bfseries\color{ink}}{\thesection}{0.7em}{}[\color{accent}\titlerule]
\titleformat{\subsection}{\large\bfseries\color{ink}}{\thesubsection}{0.7em}{}
\titlespacing*{\section}{0pt}{1.5em}{0.8em}
\titlespacing*{\subsection}{0pt}{1.2em}{0.4em}
\pagestyle{fancy}
\fancyhf{}
\fancyhead[L]{\small\color{muted}ICPC NOTEBOOK JAVA}
\fancyhead[R]{\small\color{muted}TEAM REFERENCE}
\fancyfoot[L]{\small\color{muted}Java 17}
\fancyfoot[R]{\small\color{muted}\thepage}
\renewcommand{\headrulewidth}{0.4pt}
\renewcommand{\footrulewidth}{0pt}
\lstdefinestyle{notebook}{language=Java,basicstyle=\ttfamily\footnotesize,
  keywordstyle=\bfseries\color{ink},commentstyle=\itshape\color{muted},
  stringstyle=\color{accent},backgroundcolor=\color{codebg},
  frame=single,rulecolor=\color{wash},breaklines=true,breakatwhitespace=true,
  showstringspaces=false,tabsize=4,columns=fullflexible,keepspaces=true,
  xleftmargin=0.5em,xrightmargin=0.5em,aboveskip=0.8em,belowskip=1.1em}
\lstset{style=notebook}
\newtcolorbox{guidebox}{colback=wash,colframe=wash,boxrule=0pt,
  arc=2pt,left=9pt,right=9pt,top=7pt,bottom=7pt}
\newcommand{\entrylabel}[1]{\textcolor{accent}{\bfseries #1}}
\begin{document}
\pagenumbering{roman}
\hypersetup{pageanchor=false}\begin{titlepage}\thispagestyle{empty}\centering
\vspace*{1.0in}\includegraphics[width=0.82\textwidth]{logo.png}\par
\vspace{0.8in}{\Huge\bfseries\color{ink} ICPC Notebook Java}\par
\vspace{0.25in}{\Large\color{muted} A progressive team reference for Java 17}\par
\vfill\color{accent}\rule{0.75\textwidth}{1.5pt}\par
\vspace{0.2in}\large A progressive algorithm reference for team practice\par
\vspace{0.12in}\small Profile: Intermediate \quad | \quad 31 entries\par
\end{titlepage}\hypersetup{pageanchor=true}\tableofcontents\clearpage\pagenumbering{arabic}
\section*{How to use this reference}\addcontentsline{toc}{section}{How to use this reference}
Read the algorithm's use case, cost, and caution before copying its code. The listing is the same Java source compiled by the repository tests. Use the exporter to combine the selected classes into one judge-ready Main.java.

\subsection*{Learning path: Intermediate}\addcontentsline{toc}{subsection}{Learning path}
\begin{guidebox}\entrylabel{Who this is for}\quad Teams comfortable with basics and ready to select standard algorithms under contest constraints.\par
\entrylabel{Before starting}\quad The beginner reference, asymptotic notation, sorting, BFS, and simple dynamic programming.\par
\entrylabel{Practice method}\quad Compare a fast method against a brute-force oracle on random small cases; record the first failing seed.\end{guidebox}
\textbf{Learning goals}\begin{itemize}
\item Choose between static prefix queries and update-capable trees.
\item State graph direction and weight assumptions before selecting a path algorithm.
\item Design DP states, transitions, and iteration order.
\item Use string matching and number theory with explicit numeric bounds.
\item Use offline queries, bipartite matching, and circle primitives.
\item Solve dense assignment problems and use three-dimensional vector operations.
\end{itemize}
\begin{guidebox}\textbf{Conventions}\begin{itemize}
\item Vertices and arrays are zero-indexed.
\item Ranges are half-open [left, right) unless documented otherwise.
\item Check overflow, input bounds, and recursion depth before using a snippet.
\item Choose the simplest algorithm that fits the problem constraints.
\end{itemize}\end{guidebox}
\subsection*{Choosing an algorithm}\addcontentsline{toc}{subsection}{Choosing an algorithm}
\small\begin{longtable}{>{\raggedright\arraybackslash}p{0.27\linewidth}>{\raggedright\arraybackslash}p{0.29\linewidth}>{\raggedright\arraybackslash}p{0.34\linewidth}}\toprule
\textbf{Problem signal} & \textbf{Choice} & \textbf{Typical cost}\\\midrule\endhead
Small sequence, simple condition & Direct scan or brute force & O(n) or O(n\textasciicircum{}2)\\
Frequency or membership lookup & HashMap or HashSet & Expected O(1) per operation\\
All subsets of a small set & Backtracking & O(2\textasciicircum{}n)\\
Static interval sums & Prefix sums & O(n) build, O(1) query\\
Point updates and interval sums & Fenwick or segment tree & O(log n) per operation\\
Range additions and interval sums & Lazy segment tree & O(log n) per operation\\
Unweighted graph or grid & Breadth-first search & O(V+E)\\
Incremental connectivity & Disjoint set union & Almost O(1) amortized\\
Unweighted, 0/1, nonnegative paths & BFS, 0-1 BFS, Dijkstra & O(V+E) to O((V+E) log V)\\
One pattern or many patterns & KMP, Z, Aho-Corasick & Linear in input size\\
Many static distinct range queries & Mo ordering & Typical O((n+q) sqrt n)\\
Assign left items to right items & Hopcroft-Karp matching & O(E sqrt V)\\
Dense minimum-cost assignment & Hungarian algorithm & O(n\textasciicircum{}2 m)\\
\bottomrule\end{longtable}\normalsize
\subsection*{Java contest checklist}\addcontentsline{toc}{subsection}{Java contest checklist}
\begin{guidebox}\begin{itemize}
\item Check time and memory limits before choosing object-heavy structures.
\item Use long for sums, distances, products, and counts that may exceed int.
\item Check whether input is ASCII tokens or requires Unicode handling.
\item Watch recursion depth on long paths; use an explicit stack if needed.
\item Confirm whether edges are directed, whether weights can be negative, and whether ranges are inclusive.
\item Test empty, singleton, duplicate, disconnected, and boundary cases.
\end{itemize}\end{guidebox}
\section*{Problem-solving foundations}\addcontentsline{toc}{section}{Problem-solving foundations}
\subsection*{Select by operations, not by name}
A data structure is defined by the updates and queries it supports. Prefix sums answer static sums. Fenwick trees handle point updates and prefix sums. Lazy trees handle range updates.\par
\begin{guidebox}\entrylabel{Example}\quad If queries ask for sums after changing one index, a static prefix array requires rebuilding; a Fenwick tree changes O(log n) cells.\par
\entrylabel{Try it}\quad List the update and query type before choosing a data structure.\end{guidebox}
\subsection*{Design a dynamic program}
Write the state meaning, base cases, transition, iteration order, and output location. Check whether the transition reuses a state from the current iteration.\par
\begin{guidebox}\entrylabel{Example}\quad In 0/1 knapsack, descending capacity prevents using one item twice; ascending capacity instead models unbounded use.\par
\entrylabel{Try it}\quad Derive a small table by hand and compare with the implementation.\end{guidebox}
\subsection*{Match graph assumptions to algorithms}
Unweighted edges favor BFS, 0/1 weights favor deque-based BFS, nonnegative weights permit Dijkstra, and negative weights require a different method. Directedness and disconnected vertices also matter.\par
\begin{guidebox}\entrylabel{Example}\quad A weight of -1 invalidates Dijkstra even if most edges are positive.\par
\entrylabel{Try it}\quad For each graph problem, write direction, weight range, number of vertices, and number of edges.\end{guidebox}
\clearpage\section*{Guided problems}\addcontentsline{toc}{section}{Guided problems}
\subsection*{Maintain merging teams}
\noindent\entrylabel{Task}\quad Start with vertices 0 through 3. Merge (0,1), (2,3), then (1,2). Ask whether 0 and 3 are connected.\par
\noindent\entrylabel{Result}\quad Yes\par
\noindent\entrylabel{Reasoning}\quad Each union joins two representative components. After the third union, all four vertices share one representative.\par
\noindent\entrylabel{Use}\quad DataStructures.DSU\par
\noindent\entrylabel{Cost}\quad Almost constant amortized time per operation.\par
\noindent\entrylabel{Check}\quad Compare every connectivity query with BFS on the edges inserted so far.\par
\subsection*{Use weighted shortest paths}
\noindent\entrylabel{Task}\quad Directed edges are 0 to 1 with weight 4, 0 to 2 with weight 1, and 2 to 1 with weight 1. Find the distance from 0 to 1.\par
\noindent\entrylabel{Result}\quad 2\par
\noindent\entrylabel{Reasoning}\quad The direct path costs 4. The route through vertex 2 costs 1 + 1. All weights are nonnegative, so Dijkstra applies.\par
\noindent\entrylabel{Use}\quad Graphs.dijkstra\par
\noindent\entrylabel{Cost}\quad O((V+E) log V) with a priority queue.\par
\noindent\entrylabel{Check}\quad Add a disconnected vertex and confirm it retains the unreachable sentinel.\par
\subsection*{Choose items once}
\noindent\entrylabel{Task}\quad Weights are [2,3,4], values are [4,5,6], and capacity is 5. Maximize total value with each item used at most once.\par
\noindent\entrylabel{Result}\quad 9\par
\noindent\entrylabel{Reasoning}\quad Choose the first two items for weight 5 and value 9. Update capacity downward so the current item cannot be reused.\par
\noindent\entrylabel{Use}\quad DynamicProgramming.knapsack\par
\noindent\entrylabel{Cost}\quad O(items * capacity).\par
\noindent\entrylabel{Check}\quad Enumerate all eight subsets and compare their feasible values.\par
\clearpage
\Needspace{10\baselineskip}\section{Contest setup}
\begin{guidebox}Start with a compilable single-file submission and fast input.\end{guidebox}
\Needspace{12\baselineskip}\subsection{Standalone Main template}
\begin{guidebox}\small
\entrylabel{Level}\quad Beginner\par
\entrylabel{Use when}\quad Start a Java judge submission.\par
\entrylabel{Cost}\quad Linear in the input size.\par
\entrylabel{Watch for}\quad Use the exporter to include required helper classes.\par
\end{guidebox}
\noindent\textbf{Core idea.} Start from a working input/output skeleton so effort stays on the algorithm.\par
\noindent\textbf{Worked example.} Read two integers, add them, and print one line.\par
\noindent\textbf{Practice.} Compile the template before inserting helper classes.\par
\begin{lstlisting}
import java.io.*;
import java.util.*;

/** Standalone template for judges that require a single Main.java file. */
public class Main {
    static final class FastScanner {
        private final InputStream in = new BufferedInputStream(System.in);
        private final byte[] buffer = new byte[1 << 16];
        private int pos, len;

        private int read() throws IOException {
            if (pos >= len) {
                len = in.read(buffer);
                pos = 0;
                if (len < 0) return -1;
            }
            return buffer[pos++] & 0xff;
        }

        String next() throws IOException {
            int c;
            do { c = read(); } while (c >= 0 && c <= ' ');
            if (c < 0) return null;
            StringBuilder token = new StringBuilder();
            while (c > ' ') { token.append((char) c); c = read(); }
            return token.toString();
        }

        int nextInt() throws IOException { return Integer.parseInt(next()); }
        long nextLong() throws IOException { return Long.parseLong(next()); }
    }

    public static void main(String[] args) throws Exception {
        FastScanner in = new FastScanner();
        PrintWriter out = new PrintWriter(new BufferedWriter(new OutputStreamWriter(System.out)));
        // Solve the problem here. Read with in.nextInt(), write with out.println().
        out.flush();
    }
}
\end{lstlisting}
\Needspace{12\baselineskip}\subsection{Reusable fast input}
\begin{guidebox}\small
\entrylabel{Level}\quad Beginner\par
\entrylabel{Use when}\quad Read large ASCII-token input.\par
\entrylabel{Cost}\quad Linear time; fixed-size buffer.\par
\entrylabel{Watch for}\quad Not a Unicode decoder.\par
\end{guidebox}
\noindent\textbf{Core idea.} Buffered token input reduces per-character overhead on large files.\par
\noindent\textbf{Worked example.} Read n, then process exactly n integers with nextInt().\par
\noindent\textbf{Practice.} Test EOF, negative values, and the final token without a newline.\par
\begin{lstlisting}
package notebook;

import java.io.BufferedInputStream;
import java.io.IOException;
import java.io.InputStream;

/** Fast byte input for ASCII contest tokens. Use PrintWriter/StringBuilder for output. */
public final class ContestIO {
    private final InputStream in;
    private final byte[] buffer = new byte[1 << 16];
    private int pos, len;

    public ContestIO(InputStream input) { in = new BufferedInputStream(input); }

    private int read() throws IOException {
        if (pos >= len) {
            len = in.read(buffer);
            pos = 0;
            if (len < 0) return -1;
        }
        return buffer[pos++] & 0xff;
    }

    /** Returns null at EOF. Tokens are decoded as ASCII. */
    public String next() throws IOException {
        int c;
        do { c = read(); } while (c >= 0 && c <= ' ');
        if (c < 0) return null;
        StringBuilder token = new StringBuilder();
        while (c > ' ') { token.append((char) c); c = read(); }
        return token.toString();
    }

    public int nextInt() throws IOException { return Integer.parseInt(next()); }
    public long nextLong() throws IOException { return Long.parseLong(next()); }
}
\end{lstlisting}
\Needspace{10\baselineskip}\section{Fundamentals}
\begin{guidebox}Arrays, windows, sorting, and binary search come first.\end{guidebox}
\Needspace{12\baselineskip}\subsection{Loops, arrays, prefix sums, and brute force}
\begin{guidebox}\small
\entrylabel{Level}\quad Beginner\par
\entrylabel{Use when}\quad Process a sequence or test all small pairs.\par
\entrylabel{Cost}\quad O(n) scans; O(n\textasciicircum{}2) pair search.\par
\entrylabel{Watch for}\quad A prefix array has length n+1 and answers half-open ranges.\par
\end{guidebox}
\noindent\textbf{Core idea.} A direct scan is often the clearest baseline. Prefix sums store the total before each index so one subtraction answers a range.\par
\noindent\textbf{Worked example.} For [3, -2, 5], prefix = [0, 3, 1, 6] and sum [1, 3) = 6 - 3 = 3.\par
\noindent\textbf{Practice.} Compare every range sum against a loop on small random arrays.\par
\begin{lstlisting}
package notebook;

import java.util.Arrays;

/** Small array patterns that are useful before learning advanced data structures. */
public final class BeginnerArrays {
    private BeginnerArrays() {}

    /** Sum all values. O(n) time, O(1) extra space; throws if the sum exceeds long. */
    public static long sum(int[] values) {
        long result = 0;
        for (int value : values) result = Math.addExact(result, value);
        return result;
    }

    /** Return the index of the first target, or -1. O(n) time, O(1) space. */
    public static int linearSearch(int[] values, int target) {
        for (int i = 0; i < values.length; i++) if (values[i] == target) return i;
        return -1;
    }

    /** Count occurrences of a target. O(n) time, O(1) space. */
    public static int count(int[] values, int target) {
        int result = 0;
        for (int value : values) if (value == target) result++;
        return result;
    }

    /** Reverse an array in place. O(n) time, O(1) extra space. */
    public static void reverse(int[] values) {
        for (int left = 0, right = values.length - 1; left < right; left++, right--) {
            int temp = values[left]; values[left] = values[right]; values[right] = temp;
        }
    }

    /** Inclusive prefix sums with prefix[0] = 0. O(n) time and space. */
    public static long[] prefixSums(int[] values) {
        long[] prefix = new long[values.length + 1];
        for (int i = 0; i < values.length; i++) prefix[i + 1] = Math.addExact(prefix[i], values[i]);
        return prefix;
    }

    /** Sum [left, right) using a prefix array. O(1) time. */
    public static long rangeSum(long[] prefix, int left, int right) {
        if (left < 0 || right < left || right >= prefix.length) throw new IndexOutOfBoundsException();
        return prefix[right] - prefix[left];
    }

    /** Two-sum by trying every pair; returns indices or an empty array. O(n^2) time. */
    public static int[] pairWithSum(int[] values, long target) {
        for (int i = 0; i < values.length; i++)
            for (int j = i + 1; j < values.length; j++)
                if ((long) values[i] + values[j] == target) return new int[] {i, j};
        return new int[0];
    }

    /** Copy and sort values without modifying the input. O(n log n) time. */
    public static int[] sortedCopy(int[] values) {
        int[] result = values.clone();
        Arrays.sort(result);
        return result;
    }
}
\end{lstlisting}
\Needspace{12\baselineskip}\subsection{Maps, sets, and stacks}
\begin{guidebox}\small
\entrylabel{Level}\quad Beginner\par
\entrylabel{Use when}\quad Count frequencies, remove duplicates, or check nested delimiters.\par
\entrylabel{Cost}\quad Expected O(n) time for hash collections; O(n) stack space.\par
\entrylabel{Watch for}\quad HashMap and HashSet have no guaranteed iteration order.\par
\end{guidebox}
\noindent\textbf{Core idea.} A map stores a value for each key, a set stores membership, and a stack reverses the order in which items are processed.\par
\noindent\textbf{Worked example.} For [4,2,4], the frequency map stores 4 to 2 and 2 to 1. The string ([\{\}]) is balanced.\par
\noindent\textbf{Practice.} Solve a frequency-count problem with both a map and a sorted array, then compare the tradeoffs.\par
\begin{lstlisting}
package notebook;

import java.util.ArrayDeque;
import java.util.HashMap;
import java.util.HashSet;
import java.util.Map;
import java.util.Set;

/** Everyday Java collections: map, set, and stack. */
public final class BeginnerCollections {
    private BeginnerCollections() {}

    /** Count each integer with a hash map. Expected O(n) time and O(distinct) space. */
    public static Map<Integer, Integer> frequencies(int[] values) {
        Map<Integer, Integer> result = new HashMap<>();
        for (int value : values) result.merge(value, 1, Integer::sum);
        return result;
    }

    /** Count distinct integers with a hash set. Expected O(n) time and space. */
    public static int distinctCount(int[] values) {
        Set<Integer> seen = new HashSet<>();
        for (int value : values) seen.add(value);
        return seen.size();
    }

    /** Check balanced (), [], and {} brackets. O(n) time, O(n) stack space. */
    public static boolean balancedBrackets(String text) {
        ArrayDeque<Character> stack = new ArrayDeque<>();
        for (int i = 0; i < text.length(); i++) {
            char ch = text.charAt(i);
            if (ch == '(' || ch == '[' || ch == '{') stack.push(ch);
            else if (ch == ')' || ch == ']' || ch == '}') {
                if (stack.isEmpty()) return false;
                char open = stack.pop();
                if (ch == ')' && open != '(' || ch == ']' && open != '[' || ch == '}' && open != '{')
                    return false;
            } else throw new IllegalArgumentException("expected brackets only");
        }
        return stack.isEmpty();
    }
}
\end{lstlisting}
\Needspace{12\baselineskip}\subsection{Recursion and exhaustive subsets}
\begin{guidebox}\small
\entrylabel{Level}\quad Beginner\par
\entrylabel{Use when}\quad Enumerate all choices when n is small.\par
\entrylabel{Cost}\quad O(n) factorial; O(2\textasciicircum{}n) subset search.\par
\entrylabel{Watch for}\quad Exponential search is only practical for small n; recursion uses stack space.\par
\end{guidebox}
\noindent\textbf{Core idea.} A recursive function solves a smaller instance, reaches a base case, then combines answers. For each subset element, branch into skip and take.\par
\noindent\textbf{Worked example.} For [1,2,3] and target 3, subsets [3] and [1,2] both match.\par
\noindent\textbf{Practice.} Draw the recursion tree for three values and mark its base cases.\par
\begin{lstlisting}
package notebook;

/** Recursion and exhaustive search when the input is small enough. */
public final class BeginnerBacktracking {
    private BeginnerBacktracking() {}

    /** Factorial using a smaller subproblem. O(n) time and stack; n must be in [0,20]. */
    public static long factorial(int n) {
        if (n < 0 || n > 20) throw new IllegalArgumentException("n must be in [0,20]");
        if (n <= 1) return 1;
        return n * factorial(n - 1);
    }

    /** Count subsets whose sum equals target. O(2^n) time, O(n) stack; n <= 25. */
    public static long countSubsetsWithSum(int[] values, long target) {
        if (values.length > 25) throw new IllegalArgumentException("too many subsets");
        return count(values, 0, 0, target);
    }

    private static long count(int[] values, int index, long sum, long target) {
        if (index == values.length) return sum == target ? 1 : 0;
        return count(values, index + 1, sum, target)
                + count(values, index + 1, Math.addExact(sum, values[index]), target);
    }
}
\end{lstlisting}
\Needspace{12\baselineskip}\subsection{Arrays and two pointers}
\begin{guidebox}\small
\entrylabel{Level}\quad Intermediate\par
\entrylabel{Use when}\quad Use prefix sums, differences, compression, windows, inversions, or permutations.\par
\entrylabel{Cost}\quad O(n) to O(n log n).\par
\entrylabel{Watch for}\quad Sliding-window sums require nonnegative values.\par
\end{guidebox}
\noindent\textbf{Core idea.} Prefix, window, permutation, and inversion patterns replace repeated work with maintained state.\par
\noindent\textbf{Worked example.} A nonnegative window can grow right and shrink left while its sum is too large.\par
\noindent\textbf{Practice.} State the invariant before moving either pointer.\par
\begin{lstlisting}
package notebook;

import java.util.Arrays;

/** Building blocks used before choosing a more specialized algorithm. */
public final class Fundamentals {
    private Fundamentals() {}

    /** prefix[i] is the sum of a[0..i); query [l,r) as prefix[r]-prefix[l]. O(n). */
    public static long[] prefixSums(int[] a) {
        long[] prefix = new long[a.length + 1];
        for (int i = 0; i < a.length; i++) prefix[i + 1] = prefix[i] + a[i];
        return prefix;
    }

    /** Apply [l,r) additions in O(updates+n). Each update is {l,r,delta}. */
    public static long[] rangeAdd(int n, long[][] updates) {
        long[] difference = new long[n + 1];
        for (long[] update : updates) {
            int l = Math.toIntExact(update[0]), r = Math.toIntExact(update[1]);
            if (l < 0 || l > r || r > n) throw new IndexOutOfBoundsException();
            difference[l] += update[2];
            difference[r] -= update[2];
        }
        long[] result = new long[n];
        long sum = 0;
        for (int i = 0; i < n; i++) { sum += difference[i]; result[i] = sum; }
        return result;
    }

    /** Sorted distinct values. Use Arrays.binarySearch(result, x) for compressed IDs. O(n log n). */
    public static int[] sortedUnique(int[] a) {
        int[] sorted = a.clone();
        Arrays.sort(sorted);
        int size = 0;
        for (int value : sorted)
            if (size == 0 || value != sorted[size - 1]) sorted[size++] = value;
        return Arrays.copyOf(sorted, size);
    }

    /** Two pointers: number of pairs i<j with a[i]+a[j] <= limit. Input must be sorted. O(n). */
    public static long countPairsAtMost(int[] sorted, long limit) {
        long count = 0;
        for (int l = 0, r = sorted.length - 1; l < r;) {
            if ((long) sorted[l] + sorted[r] <= limit) { count += r - l; l++; }
            else r--;
        }
        return count;
    }

    /** Longest window with sum <= limit; requires nonnegative elements. O(n). */
    public static int longestWindowAtMost(int[] a, long limit) {
        long sum = 0;
        int best = 0, left = 0;
        for (int right = 0; right < a.length; right++) {
            if (a[right] < 0) throw new IllegalArgumentException("negative element");
            sum += a[right];
            while (left <= right && sum > limit) sum -= a[left++];
            best = Math.max(best, right - left + 1);
        }
        return best;
    }

    /** Maximum nonempty subarray sum. O(n); returns 0 for empty input. */
    public static long maxSubarraySum(int[] a) {
        if (a.length == 0) return 0;
        long ending = a[0], best = ending;
        for (int i = 1; i < a.length; i++) {
            ending = Math.max(a[i], ending + a[i]);
            best = Math.max(best, ending);
        }
        return best;
    }

    /** Number of index pairs i<j with a[i]>a[j]. O(n log n). */
    public static long inversions(int[] a) {
        int[] copy = a.clone(), scratch = new int[a.length];
        return mergeCount(copy, scratch, 0, a.length);
    }

    private static long mergeCount(int[] a, int[] tmp, int lo, int hi) {
        if (hi - lo <= 1) return 0;
        int mid = (lo + hi) >>> 1;
        long count = mergeCount(a, tmp, lo, mid) + mergeCount(a, tmp, mid, hi);
        int i = lo, j = mid;
        for (int k = lo; k < hi; k++) {
            if (j == hi || (i < mid && a[i] <= a[j])) tmp[k] = a[i++];
            else { tmp[k] = a[j++]; count += mid - i; }
        }
        System.arraycopy(tmp, lo, a, lo, hi - lo);
        return count;
    }

    /** Lexicographically next permutation in place; false after wrapping to sorted order. O(n). */
    public static boolean nextPermutation(int[] a) {
        int i = a.length - 2;
        while (i >= 0 && a[i] >= a[i + 1]) i--;
        if (i >= 0) {
            int j = a.length - 1;
            while (a[j] <= a[i]) j--;
            int t = a[i]; a[i] = a[j]; a[j] = t;
        }
        for (int l = i + 1, r = a.length - 1; l < r; l++, r--) {
            int t = a[l]; a[l] = a[r]; a[r] = t;
        }
        return i >= 0;
    }

    /** Enumerate nonempty submasks of mask; mask must be nonnegative. O(2^popcount(mask)). */
    public static int[] nonemptySubmasks(int mask) {
        if (mask < 0) throw new IllegalArgumentException();
        if (Integer.bitCount(mask) > 20) throw new IllegalArgumentException("too many submasks to materialize");
        int[] result = new int[(1 << Integer.bitCount(mask)) - 1];
        int index = 0;
        for (int sub = mask; sub != 0; sub = (sub - 1) & mask) result[index++] = sub;
        return result;
    }
}
\end{lstlisting}
\Needspace{12\baselineskip}\subsection{Binary search and bounds}
\begin{guidebox}\small
\entrylabel{Level}\quad Beginner\par
\entrylabel{Use when}\quad Search sorted data or a monotone answer predicate.\par
\entrylabel{Cost}\quad O(log n) search.\par
\entrylabel{Watch for}\quad Sorted input and monotonicity are preconditions.\par
\end{guidebox}
\noindent\textbf{Core idea.} Binary search keeps an interval that is guaranteed to contain the answer.\par
\noindent\textbf{Worked example.} For sorted [1, 2, 2, 5], lowerBound(2) is 1 and upperBound(2) is 3.\par
\noindent\textbf{Practice.} Test an empty array and targets below and above every element.\par
\begin{lstlisting}
package notebook;

import java.util.function.LongPredicate;

/** Search and sorting routines. Input to searches must be sorted ascending. */
public final class SearchingSorting {
    private SearchingSorting() {}

    /** O(log n); returns any matching index, or -1. */
    public static int binarySearch(int[] a, int key) {
        int lo = 0, hi = a.length - 1;
        while (lo <= hi) {
            int mid = lo + (hi - lo) / 2;
            if (a[mid] == key) return mid;
            if (a[mid] < key) lo = mid + 1;
            else hi = mid - 1;
        }
        return -1;
    }

    /** O(log n); returns any matching index, or -1. */
    public static int ternarySearch(int[] a, int key) {
        int lo = 0, hi = a.length - 1;
        while (lo <= hi) {
            int m1 = lo + (hi - lo) / 3;
            int m2 = hi - (hi - lo) / 3;
            if (a[m1] == key) return m1;
            if (a[m2] == key) return m2;
            if (key < a[m1]) hi = m1 - 1;
            else if (key > a[m2]) lo = m2 + 1;
            else { lo = m1 + 1; hi = m2 - 1; }
        }
        return -1;
    }

    /** O(n + 65536) counting sort for Java UTF-16 char units. */
    public static void countingSort(char[] a) {
        int[] count = new int[Character.MAX_VALUE + 1];
        for (char c : a) count[c]++;
        int i = 0;
        for (int c = 0; c < count.length; c++) {
            while (count[c]-- > 0) a[i++] = (char) c;
        }
    }

    /** Lower bound: first index i with a[i] >= key. O(log n). */
    public static int lowerBound(int[] a, int key) {
        int lo = 0, hi = a.length;
        while (lo < hi) {
            int mid = lo + (hi - lo) / 2;
            if (a[mid] < key) lo = mid + 1;
            else hi = mid;
        }
        return lo;
    }

    /** First index with a[i] > key in a sorted array. O(log n). */
    public static int upperBound(int[] a, int key) {
        int lo = 0, hi = a.length;
        while (lo < hi) {
            int mid = lo + (hi - lo) / 2;
            if (a[mid] <= key) lo = mid + 1;
            else hi = mid;
        }
        return lo;
    }

    /** First x in nonnegative [low,high] satisfying a monotone false-then-true predicate. */
    public static long firstTrue(long low, long high, LongPredicate predicate) {
        if (low < 0 || low > high || !predicate.test(high)) throw new IllegalArgumentException("no true value in range");
        while (low < high) {
            long mid = low + ((high - low) >>> 1);
            if (predicate.test(mid)) high = mid;
            else low = mid + 1;
        }
        return low;
    }
}
\end{lstlisting}
\Needspace{12\baselineskip}\subsection{Interval scheduling and covering}
\begin{guidebox}\small
\entrylabel{Level}\quad Beginner\par
\entrylabel{Use when}\quad Select compatible intervals or cover intervals with few points.\par
\entrylabel{Cost}\quad O(n log n) from sorting.\par
\entrylabel{Watch for}\quad Scheduling uses half-open intervals; point covering uses closed intervals.\par
\end{guidebox}
\noindent\textbf{Core idea.} Sort by earliest finishing time. Choosing an early finish leaves as much room as possible for later intervals.\par
\noindent\textbf{Worked example.} Among [0,3), [1,2), [2,4), choose the last two.\par
\noindent\textbf{Practice.} Try to prove the choice by exchanging an optimal first interval with the greedy one.\par
\begin{lstlisting}
package notebook;

import java.util.Arrays;
import java.util.Comparator;

/** Greedy choices with simple exchange-argument proofs. */
public final class BeginnerGreedy {
    private BeginnerGreedy() {}

    /** Maximum number of non-overlapping [start,end) intervals. O(n log n) time. */
    public static int maxNonOverlapping(long[][] intervals) {
        long[][] sorted = new long[intervals.length][];
        for (int i = 0; i < intervals.length; i++) {
            if (intervals[i].length != 2 || intervals[i][0] > intervals[i][1])
                throw new IllegalArgumentException("invalid interval");
            sorted[i] = intervals[i].clone();
        }
        Arrays.sort(sorted, Comparator.comparingLong((long[] interval) -> interval[1])
                .thenComparingLong(interval -> interval[0]));
        int chosen = 0;
        long lastEnd = Long.MIN_VALUE;
        for (long[] interval : sorted) if (interval[0] >= lastEnd) {
            chosen++;
            lastEnd = interval[1];
        }
        return chosen;
    }

    /** Minimum points that cover closed intervals [left,right]. O(n log n) time. */
    public static int pointsToCover(long[][] intervals) {
        long[][] sorted = new long[intervals.length][];
        for (int i = 0; i < intervals.length; i++) {
            if (intervals[i].length != 2 || intervals[i][0] > intervals[i][1])
                throw new IllegalArgumentException("invalid interval");
            sorted[i] = intervals[i].clone();
        }
        Arrays.sort(sorted, Comparator.comparingLong(interval -> interval[1]));
        int count = 0;
        long lastPoint = 0;
        boolean hasPoint = false;
        for (long[] interval : sorted) if (!hasPoint || interval[0] > lastPoint) {
            lastPoint = interval[1];
            hasPoint = true;
            count++;
        }
        return count;
    }
}
\end{lstlisting}
\Needspace{10\baselineskip}\section{Range queries and data structures}
\begin{guidebox}Choose by update type, query type, and whether the data changes.\end{guidebox}
\Needspace{12\baselineskip}\subsection{DSU, Fenwick, and segment tree}
\begin{guidebox}\small
\entrylabel{Level}\quad Intermediate\par
\entrylabel{Use when}\quad Track connectivity, point updates, or range sums.\par
\entrylabel{Cost}\quad Almost O(1) DSU; O(log n) range operations.\par
\entrylabel{Watch for}\quad Sums use long.\par
\end{guidebox}
\noindent\textbf{Core idea.} Use DSU for merging groups, Fenwick for prefix sums with updates, and segment trees for general ranges.\par
\noindent\textbf{Worked example.} After union(0,1), find(0) and find(1) represent the same component.\par
\noindent\textbf{Practice.} Compare random updates and queries against a plain array.\par
\begin{lstlisting}
package notebook;

/** Zero-indexed contest data structures. */
public final class DataStructures {
    private DataStructures() {}

    /** Disjoint set union with path compression and union by size. */
    public static final class DSU {
        private final int[] parent, size;
        private int components;
        public DSU(int n) {
            if (n < 0) throw new IllegalArgumentException();
            parent = new int[n]; size = new int[n]; components = n;
            for (int i = 0; i < n; i++) { parent[i] = i; size[i] = 1; }
        }
        public int find(int x) {
            if (parent[x] != x) parent[x] = find(parent[x]);
            return parent[x];
        }
        public boolean union(int a, int b) {
            a = find(a); b = find(b);
            if (a == b) return false;
            if (size[a] < size[b]) { int t = a; a = b; b = t; }
            parent[b] = a; size[a] += size[b]; components--;
            return true;
        }
        public int size(int x) { return size[find(x)]; }
        public int components() { return components; }
    }

    /** Fenwick tree for point additions and half-open range sums. O(log n) per operation. */
    public static final class Fenwick {
        private final long[] tree;
        public Fenwick(int n) {
            if (n < 0) throw new IllegalArgumentException();
            tree = new long[n + 1];
        }
        public void add(int index, long delta) {
            if (index < 0 || index >= tree.length - 1) throw new IndexOutOfBoundsException();
            for (int i = index + 1; i < tree.length; i += i & -i) tree[i] += delta;
        }
        public long prefixSum(int endExclusive) {
            if (endExclusive < 0 || endExclusive >= tree.length) throw new IndexOutOfBoundsException();
            long result = 0;
            for (int i = endExclusive; i > 0; i -= i & -i) result += tree[i];
            return result;
        }
        public long rangeSum(int left, int rightExclusive) {
            if (left > rightExclusive) throw new IllegalArgumentException();
            return prefixSum(rightExclusive) - prefixSum(left);
        }
    }

    /** Iterative segment tree for point assignment and half-open range sums. */
    public static final class SegmentTree {
        private final long[] tree;
        private final int n, base;
        public SegmentTree(long[] values) {
            n = values.length;
            int power = 1;
            while (power < n) power <<= 1;
            base = power;
            tree = new long[2 * base];
            for (int i = 0; i < n; i++) tree[base + i] = values[i];
            for (int i = base - 1; i > 0; i--) tree[i] = tree[2 * i] + tree[2 * i + 1];
        }
        public void set(int index, long value) {
            if (index < 0 || index >= n) throw new IndexOutOfBoundsException();
            int p = base + index;
            tree[p] = value;
            for (p >>= 1; p > 0; p >>= 1) tree[p] = tree[2 * p] + tree[2 * p + 1];
        }
        public long sum(int left, int rightExclusive) {
            if (left < 0 || left > rightExclusive || rightExclusive > n) throw new IndexOutOfBoundsException();
            long result = 0;
            for (int l = left + base, r = rightExclusive + base; l < r; l >>= 1, r >>= 1) {
                if ((l & 1) != 0) result += tree[l++];
                if ((r & 1) != 0) result += tree[--r];
            }
            return result;
        }
    }
}
\end{lstlisting}
\Needspace{12\baselineskip}\subsection{2D prefix sums, sparse table, lazy tree}
\begin{guidebox}\small
\entrylabel{Level}\quad Intermediate\par
\entrylabel{Use when}\quad Answer rectangle sums, immutable minima, or range additions.\par
\entrylabel{Cost}\quad O(1) immutable queries; O(log n) lazy operations.\par
\entrylabel{Watch for}\quad The sparse table supports idempotent minimum.\par
\end{guidebox}
\noindent\textbf{Core idea.} Preprocessing trades build time and memory for faster repeated queries.\par
\noindent\textbf{Worked example.} A rectangle sum uses four 2D prefix values by inclusion-exclusion.\par
\noindent\textbf{Practice.} Check one-cell, whole-grid, and boundary rectangles.\par
\begin{lstlisting}
package notebook;

/** Immutable and mutable range query structures; all ranges are half-open. */
public final class RangeQueries {
    private RangeQueries() {}

    /** Rectangle sums on [top,bottom) x [left,right). O(rows*cols) build, O(1) query. */
    public static final class Prefix2D {
        private final long[][] prefix;
        public Prefix2D(int[][] grid) {
            int rows = grid.length, cols = rows == 0 ? 0 : grid[0].length;
            prefix = new long[rows + 1][cols + 1];
            for (int r = 0; r < rows; r++) {
                if (grid[r].length != cols) throw new IllegalArgumentException("ragged grid");
                for (int c = 0; c < cols; c++)
                    prefix[r + 1][c + 1] = grid[r][c] + prefix[r][c + 1]
                        + prefix[r + 1][c] - prefix[r][c];
            }
        }
        public long sum(int top, int left, int bottom, int right) {
            return prefix[bottom][right] - prefix[top][right]
                - prefix[bottom][left] + prefix[top][left];
        }
    }

    /** Static range minimum. O(n log n) build, O(1) query, O(n log n) space. */
    public static final class SparseTableMin {
        private final int[][] table;
        private final int[] log;
        public SparseTableMin(int[] a) {
            log = new int[a.length + 1];
            for (int i = 2; i < log.length; i++) log[i] = log[i / 2] + 1;
            table = new int[a.length == 0 ? 0 : log[a.length] + 1][a.length];
            if (a.length == 0) return;
            table[0] = a.clone();
            for (int k = 1; k < table.length; k++)
                for (int i = 0; i + (1 << k) <= a.length; i++)
                    table[k][i] = Math.min(table[k - 1][i], table[k - 1][i + (1 << (k - 1))]);
        }
        public int min(int left, int right) {
            if (left < 0 || left >= right || right > log.length - 1) throw new IndexOutOfBoundsException();
            int k = log[right - left];
            return Math.min(table[k][left], table[k][right - (1 << k)]);
        }
    }

    /** Range addition and range sum. O(n) build, O(log n) operations. */
    public static final class LazySegmentTree {
        private final long[] sum, lazy;
        private final int n;
        public LazySegmentTree(long[] values) {
            n = values.length;
            sum = new long[4 * Math.max(1, n)];
            lazy = new long[sum.length];
            if (n > 0) build(1, 0, n, values);
        }
        private void build(int node, int l, int r, long[] values) {
            if (r - l == 1) { sum[node] = values[l]; return; }
            int mid = (l + r) >>> 1;
            build(node * 2, l, mid, values);
            build(node * 2 + 1, mid, r, values);
            sum[node] = sum[node * 2] + sum[node * 2 + 1];
        }
        private void apply(int node, int length, long delta) {
            sum[node] += delta * length;
            lazy[node] += delta;
        }
        private void push(int node, int l, int r) {
            if (lazy[node] == 0 || r - l <= 1) return;
            int mid = (l + r) >>> 1;
            apply(node * 2, mid - l, lazy[node]);
            apply(node * 2 + 1, r - mid, lazy[node]);
            lazy[node] = 0;
        }
        public void add(int left, int right, long delta) {
            if (left < 0 || left > right || right > n) throw new IndexOutOfBoundsException();
            if (n > 0) add(1, 0, n, left, right, delta);
        }
        private void add(int node, int l, int r, int ql, int qr, long delta) {
            if (qr <= l || r <= ql) return;
            if (ql <= l && r <= qr) { apply(node, r - l, delta); return; }
            push(node, l, r);
            int mid = (l + r) >>> 1;
            add(node * 2, l, mid, ql, qr, delta);
            add(node * 2 + 1, mid, r, ql, qr, delta);
            sum[node] = sum[node * 2] + sum[node * 2 + 1];
        }
        public long sum(int left, int right) {
            if (left < 0 || left > right || right > n) throw new IndexOutOfBoundsException();
            return n == 0 ? 0 : sum(1, 0, n, left, right);
        }
        private long sum(int node, int l, int r, int ql, int qr) {
            if (qr <= l || r <= ql) return 0;
            if (ql <= l && r <= qr) return sum[node];
            push(node, l, r);
            int mid = (l + r) >>> 1;
            return sum(node * 2, l, mid, ql, qr) + sum(node * 2 + 1, mid, r, ql, qr);
        }
    }
}
\end{lstlisting}
\Needspace{12\baselineskip}\subsection{Two-dimensional Fenwick tree}
\begin{guidebox}\small
\entrylabel{Level}\quad Intermediate\par
\entrylabel{Use when}\quad Update cells and query rectangle sums.\par
\entrylabel{Cost}\quad O(log rows * log columns).\par
\entrylabel{Watch for}\quad Zero-based half-open rectangles.\par
\end{guidebox}
\noindent\textbf{Core idea.} Nested Fenwick loops aggregate a rectangle through four prefixes.\par
\noindent\textbf{Worked example.} Update (1,2), then query a rectangle containing that cell.\par
\noindent\textbf{Practice.} Compare random rectangles against a matrix.\par
\begin{lstlisting}
package notebook;

/** Point additions and rectangle sums on a fixed zero-indexed grid. */
public final class Fenwick2D {
    private final int rows, columns;
    private final long[][] bit;

    public Fenwick2D(int rows, int columns) {
        if (rows < 0 || columns < 0) throw new IllegalArgumentException();
        this.rows = rows; this.columns = columns;
        bit = new long[rows + 1][columns + 1];
    }

    /** Add delta to one cell. O(log rows * log columns). */
    public void add(int row, int column, long delta) {
        if (row < 0 || row >= rows || column < 0 || column >= columns) throw new IndexOutOfBoundsException();
        for (int r = row + 1; r <= rows; r += r & -r)
            for (int c = column + 1; c <= columns; c += c & -c)
                bit[r][c] = Math.addExact(bit[r][c], delta);
    }

    /** Sum [0,row) x [0,column). O(log rows * log columns). */
    public long prefix(int row, int column) {
        if (row < 0 || row > rows || column < 0 || column > columns) throw new IndexOutOfBoundsException();
        long result = 0;
        for (int r = row; r > 0; r -= r & -r)
            for (int c = column; c > 0; c -= c & -c)
                result = Math.addExact(result, bit[r][c]);
        return result;
    }

    /** Sum [top,bottom) x [left,right). O(log rows * log columns). */
    public long rectangle(int top, int left, int bottom, int right) {
        if (top > bottom || left > right) throw new IndexOutOfBoundsException();
        return Math.subtractExact(Math.subtractExact(prefix(bottom, right), prefix(top, right)),
                Math.subtractExact(prefix(bottom, left), prefix(top, left)));
    }
}
\end{lstlisting}
\Needspace{12\baselineskip}\subsection{Mo offline distinct queries}
\begin{guidebox}\small
\entrylabel{Level}\quad Intermediate\par
\entrylabel{Use when}\quad Answer many static range distinct queries.\par
\entrylabel{Cost}\quad Typical O((n+q) sqrt n).\par
\entrylabel{Watch for}\quad All queries must be known in advance.\par
\end{guidebox}
\noindent\textbf{Core idea.} Order queries by left block to limit pointer movement.\par
\noindent\textbf{Worked example.} For [1,2,1], distinct [0,3) is 2.\par
\noindent\textbf{Practice.} Compare with a HashSet on small arrays.\par
\begin{lstlisting}
package notebook;

import java.util.Arrays;
import java.util.HashMap;
import java.util.Map;

/** Offline static range queries ordered by Mo's square-root decomposition. */
public final class MoQueries {
    private MoQueries() {}

    /** Distinct-value counts for half-open [left,right) queries. O((n+q)*sqrt(n)) typical. */
    public static int[] distinctCounts(int[] values, int[][] ranges) {
        int n = values.length, block = Math.max(1, (int) Math.sqrt(n));
        Map<Integer, Integer> compressed = new HashMap<>();
        int[] ids = new int[n];
        for (int i = 0; i < n; i++) {
            if (!compressed.containsKey(values[i])) compressed.put(values[i], compressed.size());
            ids[i] = compressed.get(values[i]);
        }
        Integer[] order = new Integer[ranges.length];
        for (int i = 0; i < ranges.length; i++) {
            if (ranges[i].length != 2 || ranges[i][0] < 0 || ranges[i][0] > ranges[i][1]
                    || ranges[i][1] > n) throw new IllegalArgumentException("invalid range");
            order[i] = i;
        }
        Arrays.sort(order, (a, b) -> {
            int leftBlock = ranges[a][0] / block, rightBlock = ranges[b][0] / block;
            if (leftBlock != rightBlock) return Integer.compare(leftBlock, rightBlock);
            return leftBlock % 2 == 0 ? Integer.compare(ranges[a][1], ranges[b][1])
                    : Integer.compare(ranges[b][1], ranges[a][1]);
        });
        int[] frequency = new int[compressed.size()], result = new int[ranges.length];
        int left = 0, right = 0, distinct = 0;
        for (int index : order) {
            int targetLeft = ranges[index][0], targetRight = ranges[index][1];
            while (left > targetLeft) if (frequency[ids[--left]]++ == 0) distinct++;
            while (right < targetRight) if (frequency[ids[right++]]++ == 0) distinct++;
            while (left < targetLeft) if (--frequency[ids[left++]] == 0) distinct--;
            while (right > targetRight) if (--frequency[ids[--right]] == 0) distinct--;
            result[index] = distinct;
        }
        return result;
    }
}
\end{lstlisting}
\Needspace{10\baselineskip}\section{Dynamic programming}
\begin{guidebox}State, transition, base case, and iteration order should be explicit.\end{guidebox}
\Needspace{12\baselineskip}\subsection{First dynamic programs}
\begin{guidebox}\small
\entrylabel{Level}\quad Beginner\par
\entrylabel{Use when}\quad Reuse answers to smaller states for sequences, coin change, or grids.\par
\entrylabel{Cost}\quad O(n), O(amount * coin types), or O(rows * columns).\par
\entrylabel{Watch for}\quad Path counts can overflow long; coin denominations must be positive.\par
\end{guidebox}
\noindent\textbf{Core idea.} Define the state, base case, transition, and answer before writing loops. An iterative table avoids repeated recursion.\par
\noindent\textbf{Worked example.} For coins 1, 3, 4, amount 6 needs two coins: 3 + 3.\par
\noindent\textbf{Practice.} Write a table by hand for a small input and compare it with the program.\par
\begin{lstlisting}
package notebook;

import java.util.Arrays;

/** First DP examples: identify a state, base case, transition, and answer. */
public final class BeginnerDynamicProgramming {
    private BeginnerDynamicProgramming() {}

    /** Fibonacci F(n), with F(0)=0 and F(1)=1. O(n) time; n must be at most 92. */
    public static long fibonacci(int n) {
        if (n < 0 || n > 92) throw new IllegalArgumentException("n must be in [0,92]");
        long previous = 0, current = 1;
        for (int i = 0; i < n; i++) { long next = previous + current; previous = current; current = next; }
        return previous;
    }

    /** Minimum coins for an amount, or -1 if impossible. O(amount * coins.length) time. */
    public static int minimumCoins(int[] coins, int amount) {
        if (amount < 0) throw new IllegalArgumentException("negative amount");
        int[] dp = new int[amount + 1];
        Arrays.fill(dp, amount + 1);
        dp[0] = 0;
        for (int value = 1; value <= amount; value++)
            for (int coin : coins) {
                if (coin <= 0) throw new IllegalArgumentException("coins must be positive");
                if (coin <= value) dp[value] = Math.min(dp[value], dp[value - coin] + 1);
            }
        return dp[amount] > amount ? -1 : dp[amount];
    }

    /** Number of paths in a grid, moving right or down around blocked cells. O(rows*columns). */
    public static long gridPaths(boolean[][] blocked) {
        int rows = blocked.length;
        if (rows == 0) return 0;
        int columns = blocked[0].length;
        for (boolean[] row : blocked) if (row.length != columns) throw new IllegalArgumentException("ragged grid");
        if (columns == 0) return 0;
        long[] dp = new long[columns];
        dp[0] = blocked[0][0] ? 0 : 1;
        for (int row = 0; row < rows; row++)
            for (int column = 0; column < columns; column++) {
                if (blocked[row][column]) dp[column] = 0;
                else if (column > 0) dp[column] = Math.addExact(dp[column], dp[column - 1]);
            }
        return dp[columns - 1];
    }
}
\end{lstlisting}
\Needspace{12\baselineskip}\subsection{Classic DP patterns}
\begin{guidebox}\small
\entrylabel{Level}\quad Intermediate\par
\entrylabel{Use when}\quad Solve knapsack, subsequence, edit, partition, and counting problems.\par
\entrylabel{Cost}\quad Often O(n*capacity) or O(n\textasciicircum{}2).\par
\entrylabel{Watch for}\quad Counts can exceed long.\par
\end{guidebox}
\noindent\textbf{Core idea.} Many contest DPs differ only in the state and transition; order must respect dependencies.\par
\noindent\textbf{Worked example.} In 0/1 knapsack, iterate capacity downward so an item is used at most once.\par
\noindent\textbf{Practice.} Write a brute-force subset oracle for small n.\par
\begin{lstlisting}
package notebook;

/** Small contest DP routines. Counts are long where practical and may overflow. */
public final class DynamicProgramming {
    private DynamicProgramming() {}

    /** Number of unordered combinations; positive coin values required. O(coins * target). */
    public static long coinChange(int[] coins, int target) {
        if (target < 0) return 0;
        long[] dp = new long[target + 1];
        dp[0] = 1;
        for (int coin : coins) {
            if (coin <= 0) throw new IllegalArgumentException("coin must be positive");
            for (int s = coin; s <= target; s++) dp[s] += dp[s - coin];
        }
        return dp[target];
    }

    /** Number of ordered ways to climb using steps of 1, 2, or 3. O(n). */
    public static long countWays(int n) {
        if (n < 0) return 0;
        long[] dp = new long[n + 1];
        dp[0] = 1;
        for (int i = 1; i <= n; i++)
            for (int step = 1; step <= 3; step++)
                if (i >= step) dp[i] += dp[i - step];
        return dp[n];
    }

    /** Levenshtein distance. O(a.length * b.length) time, O(b.length) space. */
    public static int editDistance(String a, String b) {
        int[] prev = new int[b.length() + 1];
        for (int j = 0; j <= b.length(); j++) prev[j] = j;
        for (int i = 1; i <= a.length(); i++) {
            int[] cur = new int[b.length() + 1];
            cur[0] = i;
            for (int j = 1; j <= b.length(); j++)
                cur[j] = a.charAt(i - 1) == b.charAt(j - 1) ? prev[j - 1]
                    : 1 + Math.min(prev[j - 1], Math.min(prev[j], cur[j - 1]));
            prev = cur;
        }
        return prev[b.length()];
    }

    /** 0/1 knapsack with nonnegative weights and capacity. O(items * capacity). */
    public static long knapsack(int capacity, int[] weights, int[] values) {
        if (capacity < 0 || weights.length != values.length) throw new IllegalArgumentException();
        long[] dp = new long[capacity + 1];
        for (int i = 0; i < weights.length; i++) {
            if (weights[i] < 0) throw new IllegalArgumentException("negative weight");
            for (int w = capacity; w >= weights[i]; w--)
                dp[w] = Math.max(dp[w], dp[w - weights[i]] + values[i]);
        }
        return dp[capacity];
    }

    /** Longest common subsequence length. O(mn) time, O(n) space. */
    public static int lcs(String a, String b) {
        int[] dp = new int[b.length() + 1];
        for (int i = 1; i <= a.length(); i++) {
            int diag = 0;
            for (int j = 1; j <= b.length(); j++) {
                int old = dp[j];
                dp[j] = a.charAt(i - 1) == b.charAt(j - 1) ? diag + 1
                    : Math.max(dp[j], dp[j - 1]);
                diag = old;
            }
        }
        return dp[b.length()];
    }

    /** Strictly increasing subsequence length. O(n log n). */
    public static int lis(int[] a) {
        int[] tails = new int[a.length];
        int size = 0;
        for (int x : a) {
            int lo = 0, hi = size;
            while (lo < hi) {
                int mid = (lo + hi) >>> 1;
                if (tails[mid] < x) lo = mid + 1;
                else hi = mid;
            }
            tails[lo] = x;
            if (lo == size) size++;
        }
        return size;
    }

    /** Maximum product after at least one cut; n >= 2. O(n). */
    public static long maxProductCutting(int n) {
        if (n < 2) throw new IllegalArgumentException();
        if (n <= 3) return n - 1;
        long result = 1;
        while (n > 4) { result = Math.multiplyExact(result, 3); n -= 3; }
        return Math.multiplyExact(result, n);
    }

    /** Minimum difference between subset sums; nonnegative values. O(n * sum) time. */
    public static int minimumPartition(int[] a) {
        int sum = 0;
        for (int x : a) {
            if (x < 0) throw new IllegalArgumentException("negative value");
            sum = Math.addExact(sum, x);
        }
        boolean[] dp = new boolean[sum / 2 + 1];
        dp[0] = true;
        for (int x : a)
            for (int s = dp.length - 1; s >= x; s--) dp[s] |= dp[s - x];
        for (int s = dp.length - 1; s >= 0; s--) if (dp[s]) return sum - 2 * s;
        throw new AssertionError();
    }

    /** Subset sum for nonnegative values and target. O(n * target). */
    public static boolean subsetSum(int[] a, int target) {
        if (target < 0) return false;
        boolean[] dp = new boolean[target + 1];
        dp[0] = true;
        for (int x : a) {
            if (x < 0) throw new IllegalArgumentException("negative value");
            for (int s = target; s >= x; s--) dp[s] |= dp[s - x];
        }
        return dp[target];
    }
}
\end{lstlisting}
\Needspace{10\baselineskip}\section{Graphs}
\begin{guidebox}Distinguish directed and undirected edges before choosing an algorithm.\end{guidebox}
\Needspace{12\baselineskip}\subsection{Adjacency lists, BFS, and grid flood fill}
\begin{guidebox}\small
\entrylabel{Level}\quad Beginner\par
\entrylabel{Use when}\quad Explore an unweighted graph or count regions in a grid.\par
\entrylabel{Cost}\quad O(V+E) for graphs; O(rows * columns) for grids.\par
\entrylabel{Watch for}\quad For undirected graphs, insert both directions; BFS distance counts edges.\par
\end{guidebox}
\noindent\textbf{Core idea.} An adjacency list stores each vertex's neighbors. BFS visits vertices in layers, so its first visit gives the shortest unweighted distance.\par
\noindent\textbf{Worked example.} For edges 0-1 and 1-2, BFS from 0 returns distances 0, 1, 2.\par
\noindent\textbf{Practice.} Draw a graph with an unreachable vertex and check the returned -1 distance.\par
\begin{lstlisting}
package notebook;

import java.util.ArrayDeque;
import java.util.ArrayList;
import java.util.Arrays;
import java.util.List;

/** First graph routines using adjacency lists and an explicit queue. */
public final class BeginnerGraphs {
    private BeginnerGraphs() {}

    /** Build an undirected adjacency list from zero-indexed edges. O(V+E) time and space. */
    public static List<List<Integer>> undirectedGraph(int vertices, int[][] edges) {
        if (vertices < 0) throw new IllegalArgumentException("negative vertex count");
        List<List<Integer>> graph = new ArrayList<>();
        for (int i = 0; i < vertices; i++) graph.add(new ArrayList<>());
        for (int[] edge : edges) {
            if (edge.length != 2) throw new IllegalArgumentException("edge must contain two vertices");
            graph.get(edge[0]).add(edge[1]);
            graph.get(edge[1]).add(edge[0]);
        }
        return graph;
    }

    /** BFS distances from source; -1 means unreachable. O(V+E) time. */
    public static int[] distances(List<List<Integer>> graph, int source) {
        int[] distance = new int[graph.size()];
        Arrays.fill(distance, -1);
        ArrayDeque<Integer> queue = new ArrayDeque<>();
        distance[source] = 0;
        queue.add(source);
        while (!queue.isEmpty()) {
            int vertex = queue.remove();
            for (int next : graph.get(vertex)) if (distance[next] == -1) {
                distance[next] = distance[vertex] + 1;
                queue.add(next);
            }
        }
        return distance;
    }

    /** Count four-neighbor components of true cells. O(rows*columns) time and space. */
    public static int gridComponents(boolean[][] open) {
        int rows = open.length;
        if (rows == 0) return 0;
        int columns = open[0].length;
        for (boolean[] row : open) if (row.length != columns) throw new IllegalArgumentException("ragged grid");
        boolean[][] seen = new boolean[rows][columns];
        int[] dr = {-1, 1, 0, 0}, dc = {0, 0, -1, 1};
        ArrayDeque<int[]> queue = new ArrayDeque<>();
        int components = 0;
        for (int row = 0; row < rows; row++) for (int column = 0; column < columns; column++) {
            if (!open[row][column] || seen[row][column]) continue;
            components++;
            seen[row][column] = true;
            queue.add(new int[] {row, column});
            while (!queue.isEmpty()) {
                int[] cell = queue.remove();
                for (int direction = 0; direction < 4; direction++) {
                    int nextRow = cell[0] + dr[direction], nextColumn = cell[1] + dc[direction];
                    if (nextRow < 0 || nextRow >= rows || nextColumn < 0 || nextColumn >= columns
                            || !open[nextRow][nextColumn] || seen[nextRow][nextColumn]) continue;
                    seen[nextRow][nextColumn] = true;
                    queue.add(new int[] {nextRow, nextColumn});
                }
            }
        }
        return components;
    }
}
\end{lstlisting}
\Needspace{12\baselineskip}\subsection{Traversal, shortest paths, MST}
\begin{guidebox}\small
\entrylabel{Level}\quad Intermediate\par
\entrylabel{Use when}\quad Use BFS, DFS, topological order, Dijkstra, or Kruskal.\par
\entrylabel{Cost}\quad O(V+E) to O(E log V).\par
\entrylabel{Watch for}\quad Dijkstra requires nonnegative weights.\par
\entrylabel{Dependencies}\quad structures\par
\end{guidebox}
\noindent\textbf{Core idea.} Represent edges explicitly, then select traversal or optimization from the weight rules.\par
\noindent\textbf{Worked example.} BFS finds unweighted distances; Dijkstra extends this to nonnegative weights.\par
\noindent\textbf{Practice.} Test disconnected vertices and parallel edges.\par
\begin{lstlisting}
package notebook;

import java.util.ArrayDeque;
import java.util.ArrayList;
import java.util.Arrays;
import java.util.Comparator;
import java.util.List;
import java.util.PriorityQueue;

/** Zero-indexed graph algorithms. */
public final class Graphs {
    private Graphs() {}

    /** Directed adjacency list. For undirected graphs add both directions. */
    public static List<List<Integer>> newGraph(int n) {
        if (n < 0) throw new IllegalArgumentException();
        List<List<Integer>> graph = new ArrayList<>();
        for (int i = 0; i < n; i++) graph.add(new ArrayList<>());
        return graph;
    }

    /** Unweighted shortest-path distances; unreachable vertices have -1. O(V+E). */
    public static int[] bfs(List<List<Integer>> graph, int start) {
        int[] dist = new int[graph.size()];
        Arrays.fill(dist, -1);
        ArrayDeque<Integer> queue = new ArrayDeque<>();
        dist[start] = 0; queue.add(start);
        while (!queue.isEmpty()) {
            int v = queue.remove();
            for (int to : graph.get(v)) if (dist[to] == -1) {
                dist[to] = dist[v] + 1; queue.add(to);
            }
        }
        return dist;
    }

    /** Iterative DFS visitation order. O(V+E). */
    public static List<Integer> dfs(List<List<Integer>> graph, int start) {
        boolean[] seen = new boolean[graph.size()];
        ArrayDeque<Integer> stack = new ArrayDeque<>();
        List<Integer> order = new ArrayList<>();
        stack.push(start);
        while (!stack.isEmpty()) {
            int v = stack.pop();
            if (seen[v]) continue;
            seen[v] = true; order.add(v);
            List<Integer> neighbors = graph.get(v);
            for (int i = neighbors.size() - 1; i >= 0; i--) if (!seen[neighbors.get(i)]) stack.push(neighbors.get(i));
        }
        return order;
    }

    /** Kahn topological order; throws if the directed graph contains a cycle. */
    public static List<Integer> topologicalSort(List<List<Integer>> graph) {
        int[] indegree = new int[graph.size()];
        for (List<Integer> edges : graph) for (int to : edges) indegree[to]++;
        ArrayDeque<Integer> queue = new ArrayDeque<>();
        for (int i = 0; i < indegree.length; i++) if (indegree[i] == 0) queue.add(i);
        List<Integer> result = new ArrayList<>();
        while (!queue.isEmpty()) {
            int v = queue.remove(); result.add(v);
            for (int to : graph.get(v)) if (--indegree[to] == 0) queue.add(to);
        }
        if (result.size() != graph.size()) throw new IllegalArgumentException("graph has a cycle");
        return result;
    }

    public record WeightedEdge(int to, long weight) {}

    /** Dijkstra; requires nonnegative weights. Unreachable vertices have Long.MAX_VALUE. */
    public static long[] dijkstra(List<List<WeightedEdge>> graph, int start) {
        long[] dist = new long[graph.size()];
        Arrays.fill(dist, Long.MAX_VALUE);
        PriorityQueue<long[]> queue = new PriorityQueue<>(Comparator.comparingLong(a -> a[0]));
        dist[start] = 0; queue.add(new long[] {0, start});
        while (!queue.isEmpty()) {
            long[] current = queue.remove();
            long d = current[0]; int v = (int) current[1];
            if (d != dist[v]) continue;
            for (WeightedEdge edge : graph.get(v)) {
                if (edge.weight < 0) throw new IllegalArgumentException("negative edge");
                if (d > Long.MAX_VALUE - edge.weight) continue;
                long candidate = d + edge.weight;
                if (candidate < dist[edge.to]) {
                    dist[edge.to] = candidate;
                    queue.add(new long[] {candidate, edge.to});
                }
            }
        }
        return dist;
    }

    public record Edge(int from, int to, long weight) {}

    /** Kruskal MST; throws if disconnected. O(E log E). */
    public static long kruskal(int n, List<Edge> edges) {
        List<Edge> sorted = new ArrayList<>(edges);
        sorted.sort(Comparator.comparingLong(Edge::weight));
        DataStructures.DSU dsu = new DataStructures.DSU(n);
        long total = 0;
        for (Edge edge : sorted) if (dsu.union(edge.from, edge.to)) total = Math.addExact(total, edge.weight);
        if (n > 0 && dsu.components() != 1) throw new IllegalArgumentException("disconnected graph");
        return total;
    }
}
\end{lstlisting}
\Needspace{12\baselineskip}\subsection{Components, bipartite, 0-1 BFS, Floyd, Euler}
\begin{guidebox}\small
\entrylabel{Level}\quad Intermediate\par
\entrylabel{Use when}\quad Handle graph properties and special shortest paths.\par
\entrylabel{Cost}\quad Usually O(V+E); Floyd is O(V\textasciicircum{}3).\par
\entrylabel{Watch for}\quad Undirected edges require both directions.\par
\entrylabel{Dependencies}\quad graphs\par
\end{guidebox}
\noindent\textbf{Core idea.} Special graph structure can make a hard general problem simpler.\par
\noindent\textbf{Worked example.} A deque handles 0-weight edges at the front and 1-weight edges at the back.\par
\noindent\textbf{Practice.} Compare 0-1 BFS against Dijkstra on small graphs.\par
\begin{lstlisting}
package notebook;

import java.util.ArrayDeque;
import java.util.ArrayList;
import java.util.Arrays;
import java.util.Collections;
import java.util.List;

/** Graph routines that complement BFS, DFS, Dijkstra, and Kruskal. */
public final class GraphEssentials {
    private GraphEssentials() {}

    /** Component labels for an undirected adjacency list. O(V+E). */
    public static int[] components(List<List<Integer>> graph) {
        int[] label = new int[graph.size()];
        Arrays.fill(label, -1);
        int count = 0;
        ArrayDeque<Integer> queue = new ArrayDeque<>();
        for (int start = 0; start < graph.size(); start++) {
            if (label[start] != -1) continue;
            label[start] = count; queue.add(start);
            while (!queue.isEmpty()) {
                int v = queue.remove();
                for (int to : graph.get(v)) if (label[to] == -1) {
                    label[to] = count; queue.add(to);
                }
            }
            count++;
        }
        return label;
    }

    /** Two-color an undirected graph; return null if it is not bipartite. O(V+E). */
    public static int[] bipartiteColors(List<List<Integer>> graph) {
        int[] color = new int[graph.size()];
        Arrays.fill(color, -1);
        ArrayDeque<Integer> queue = new ArrayDeque<>();
        for (int start = 0; start < graph.size(); start++) {
            if (color[start] != -1) continue;
            color[start] = 0; queue.add(start);
            while (!queue.isEmpty()) {
                int v = queue.remove();
                for (int to : graph.get(v)) {
                    if (color[to] == color[v]) return null;
                    if (color[to] == -1) { color[to] = color[v] ^ 1; queue.add(to); }
                }
            }
        }
        return color;
    }

    public record BinaryEdge(int to, int weight) {}

    /** Shortest paths with edge weights exactly 0 or 1. O(V+E). */
    public static int[] zeroOneBfs(List<List<BinaryEdge>> graph, int source) {
        int[] dist = new int[graph.size()];
        Arrays.fill(dist, Integer.MAX_VALUE);
        ArrayDeque<Integer> deque = new ArrayDeque<>();
        dist[source] = 0; deque.add(source);
        while (!deque.isEmpty()) {
            int v = deque.removeFirst();
            for (BinaryEdge edge : graph.get(v)) {
                if (edge.weight != 0 && edge.weight != 1) throw new IllegalArgumentException("weight must be 0 or 1");
                int candidate = dist[v] + edge.weight;
                if (candidate < dist[edge.to]) {
                    dist[edge.to] = candidate;
                    if (edge.weight == 0) deque.addFirst(edge.to);
                    else deque.addLast(edge.to);
                }
            }
        }
        return dist;
    }

    /** All-pairs shortest paths; Long.MAX_VALUE means unreachable. O(V^3). */
    public static long[][] floydWarshall(long[][] weights) {
        int n = weights.length;
        long[][] dist = new long[n][n];
        for (int i = 0; i < n; i++) {
            if (weights[i].length != n) throw new IllegalArgumentException("matrix must be square");
            dist[i] = weights[i].clone();
            dist[i][i] = Math.min(dist[i][i], 0);
        }
        for (int k = 0; k < n; k++)
            for (int i = 0; i < n; i++)
                for (int j = 0; j < n; j++) {
                    if (dist[i][k] == Long.MAX_VALUE || dist[k][j] == Long.MAX_VALUE) continue;
                    long candidate = Math.addExact(dist[i][k], dist[k][j]);
                    if (candidate < dist[i][j]) dist[i][j] = candidate;
                }
        for (int i = 0; i < n; i++) if (dist[i][i] < 0) throw new IllegalArgumentException("negative cycle");
        return dist;
    }

    /** Eulerian trail in a directed multigraph, or empty list if none exists. O(V+E). */
    public static List<Integer> directedEulerTrail(int n, int[][] edges) {
        if (n == 0) return List.of();
        List<List<Integer>> graph = Graphs.newGraph(n);
        int[] in = new int[n], out = new int[n], next = new int[n];
        for (int[] edge : edges) {
            graph.get(edge[0]).add(edge[1]); out[edge[0]]++; in[edge[1]]++;
        }
        int start = -1, end = -1;
        for (int v = 0; v < n; v++) {
            int delta = out[v] - in[v];
            if (delta == 1) { if (start != -1) return List.of(); start = v; }
            else if (delta == -1) { if (end != -1) return List.of(); end = v; }
            else if (delta != 0) return List.of();
        }
        if ((start == -1) != (end == -1)) return List.of();
        if (start == -1) for (int v = 0; v < n; v++) if (out[v] > 0) { start = v; break; }
        if (start == -1) return List.of(0);
        ArrayDeque<Integer> stack = new ArrayDeque<>();
        List<Integer> trail = new ArrayList<>();
        stack.push(start);
        while (!stack.isEmpty()) {
            int v = stack.peek();
            if (next[v] < graph.get(v).size()) stack.push(graph.get(v).get(next[v]++));
            else trail.add(stack.pop());
        }
        if (trail.size() != edges.length + 1) return List.of();
        Collections.reverse(trail);
        return trail;
    }
}
\end{lstlisting}
\Needspace{12\baselineskip}\subsection{Bipartite matching and vertex cover}
\begin{guidebox}\small
\entrylabel{Level}\quad Intermediate\par
\entrylabel{Use when}\quad Assign left vertices to right vertices.\par
\entrylabel{Cost}\quad O(E sqrt V) matching.\par
\entrylabel{Watch for}\quad The DFS phase is recursive.\par
\end{guidebox}
\noindent\textbf{Core idea.} Layered augmenting paths build a maximum matching; alternating reachability gives a minimum cover.\par
\noindent\textbf{Worked example.} Edges 0-a, 1-a, 1-b allow a matching of size two.\par
\noindent\textbf{Practice.} Brute-force tiny matchings and validate the cover.\par
\begin{lstlisting}
package notebook;

import java.util.ArrayDeque;
import java.util.ArrayList;
import java.util.Arrays;
import java.util.List;

/** Hopcroft-Karp matching and Konig minimum vertex cover in bipartite graphs. */
public final class BipartiteMatching {
    public record Cover(int[] left, int[] right) {}
    private final List<List<Integer>> edges;
    private final int[] leftMatch, rightMatch, distance;
    private int endDistance;

    public BipartiteMatching(int leftSize, int rightSize) {
        if (leftSize < 0 || rightSize < 0) throw new IllegalArgumentException();
        edges = new ArrayList<>();
        for (int i = 0; i < leftSize; i++) edges.add(new ArrayList<>());
        leftMatch = new int[leftSize]; rightMatch = new int[rightSize]; distance = new int[leftSize];
        Arrays.fill(leftMatch, -1); Arrays.fill(rightMatch, -1);
    }

    public void addEdge(int left, int right) {
        if (right < 0 || right >= rightMatch.length) throw new IndexOutOfBoundsException();
        edges.get(left).add(right);
    }

    /** Maximum matching size. O(E*sqrt(V)); DFS is recursive. */
    public int maximumMatching() {
        Arrays.fill(leftMatch, -1); Arrays.fill(rightMatch, -1);
        int size = 0;
        while (layers())
            for (int left = 0; left < leftMatch.length; left++)
                if (leftMatch[left] == -1 && augment(left)) size++;
        return size;
    }

    private boolean layers() {
        Arrays.fill(distance, -1);
        ArrayDeque<Integer> queue = new ArrayDeque<>();
        for (int left = 0; left < leftMatch.length; left++) if (leftMatch[left] == -1) {
            distance[left] = 0; queue.add(left);
        }
        endDistance = Integer.MAX_VALUE;
        while (!queue.isEmpty()) {
            int left = queue.remove();
            if (distance[left] >= endDistance) continue;
            for (int right : edges.get(left)) {
                int partner = rightMatch[right];
                if (partner == -1) endDistance = distance[left] + 1;
                else if (distance[partner] == -1) {
                    distance[partner] = distance[left] + 1;
                    queue.add(partner);
                }
            }
        }
        return endDistance != Integer.MAX_VALUE;
    }

    private boolean augment(int left) {
        for (int right : edges.get(left)) {
            int partner = rightMatch[right];
            if (partner == -1 && distance[left] + 1 == endDistance
                    || partner != -1 && distance[partner] == distance[left] + 1 && augment(partner)) {
                leftMatch[left] = right; rightMatch[right] = left; return true;
            }
        }
        distance[left] = -1;
        return false;
    }

    /** Minimum vertex cover; calls maximumMatching first. O(E*sqrt(V)). */
    public Cover minimumVertexCover() {
        maximumMatching();
        boolean[] seenLeft = new boolean[leftMatch.length], seenRight = new boolean[rightMatch.length];
        ArrayDeque<Integer> queue = new ArrayDeque<>();
        for (int left = 0; left < leftMatch.length; left++) if (leftMatch[left] == -1) {
            seenLeft[left] = true; queue.add(left);
        }
        while (!queue.isEmpty()) {
            int left = queue.remove();
            for (int right : edges.get(left)) {
                if (leftMatch[left] == right || seenRight[right]) continue;
                seenRight[right] = true;
                int partner = rightMatch[right];
                if (partner != -1 && !seenLeft[partner]) { seenLeft[partner] = true; queue.add(partner); }
            }
        }
        int leftCount = 0, rightCount = 0;
        for (boolean seen : seenLeft) if (!seen) leftCount++;
        for (boolean seen : seenRight) if (seen) rightCount++;
        int[] coverLeft = new int[leftCount], coverRight = new int[rightCount];
        leftCount = rightCount = 0;
        for (int i = 0; i < seenLeft.length; i++) if (!seenLeft[i]) coverLeft[leftCount++] = i;
        for (int i = 0; i < seenRight.length; i++) if (seenRight[i]) coverRight[rightCount++] = i;
        return new Cover(coverLeft, coverRight);
    }

    public int[] leftMatches() { return leftMatch.clone(); }
}
\end{lstlisting}
\Needspace{12\baselineskip}\subsection{Minimum-cost assignment}
\begin{guidebox}\small
\entrylabel{Level}\quad Intermediate\par
\entrylabel{Use when}\quad Assign every worker to a distinct task at minimum total cost.\par
\entrylabel{Cost}\quad O(rows\textasciicircum{}2 * columns) for rows <= columns.\par
\entrylabel{Watch for}\quad Dense matrix model; arithmetic overflow is reported.\par
\end{guidebox}
\noindent\textbf{Core idea.} Dual potentials maintain nonnegative reduced costs while augmenting the assignment.\par
\noindent\textbf{Worked example.} For [[4,1,3],[2,0,5]], the minimum assignment cost is 3.\par
\noindent\textbf{Practice.} Compare all column permutations for matrices with at most seven rows.\par
\begin{lstlisting}
package notebook;

import java.util.Arrays;

/** Minimum-cost assignment of every row to a distinct column. */
public final class HungarianAlgorithm {
    private HungarianAlgorithm() {}
    public record Assignment(long cost, int[] columnForRow) {}

    /** Rectangular dense assignment, rows <= columns. O(rows^2 * columns). */
    public static Assignment minimize(long[][] costs) {
        int rows = costs.length;
        if (rows == 0) return new Assignment(0, new int[0]);
        int columns = costs[0].length;
        if (columns < rows) throw new IllegalArgumentException("rows must not exceed columns");
        for (long[] row : costs) if (row.length != columns) throw new IllegalArgumentException("ragged cost matrix");
        long[] rowPotential = new long[rows + 1], columnPotential = new long[columns + 1];
        int[] matchedRow = new int[columns + 1], previousColumn = new int[columns + 1];
        for (int row = 1; row <= rows; row++) {
            matchedRow[0] = row;
            long[] minimum = new long[columns + 1];
            Arrays.fill(minimum, Long.MAX_VALUE);
            boolean[] used = new boolean[columns + 1];
            int currentColumn = 0;
            do {
                used[currentColumn] = true;
                int currentRow = matchedRow[currentColumn], nextColumn = -1;
                long delta = Long.MAX_VALUE;
                for (int column = 1; column <= columns; column++) if (!used[column]) {
                    long reduced = Math.subtractExact(Math.subtractExact(costs[currentRow - 1][column - 1],
                            rowPotential[currentRow]), columnPotential[column]);
                    if (reduced < minimum[column]) {
                        minimum[column] = reduced;
                        previousColumn[column] = currentColumn;
                    }
                    if (minimum[column] < delta) { delta = minimum[column]; nextColumn = column; }
                }
                if (nextColumn < 0 || delta == Long.MAX_VALUE) throw new IllegalArgumentException("no assignment");
                for (int column = 0; column <= columns; column++) {
                    if (used[column]) {
                        rowPotential[matchedRow[column]] = Math.addExact(rowPotential[matchedRow[column]], delta);
                        columnPotential[column] = Math.subtractExact(columnPotential[column], delta);
                    } else if (minimum[column] != Long.MAX_VALUE)
                        minimum[column] = Math.subtractExact(minimum[column], delta);
                }
                currentColumn = nextColumn;
            } while (matchedRow[currentColumn] != 0);
            do {
                int previous = previousColumn[currentColumn];
                matchedRow[currentColumn] = matchedRow[previous];
                currentColumn = previous;
            } while (currentColumn != 0);
        }
        int[] assignment = new int[rows];
        for (int column = 1; column <= columns; column++)
            if (matchedRow[column] != 0) assignment[matchedRow[column] - 1] = column - 1;
        long total = 0;
        for (int row = 0; row < rows; row++) total = Math.addExact(total, costs[row][assignment[row]]);
        return new Assignment(total, assignment);
    }
}
\end{lstlisting}
\Needspace{10\baselineskip}\section{Strings}
\begin{guidebox}Choose a one-pattern, many-pattern, palindrome, or suffix method.\end{guidebox}
\Needspace{12\baselineskip}\subsection{Characters, frequencies, and direct matching}
\begin{guidebox}\small
\entrylabel{Level}\quad Beginner\par
\entrylabel{Use when}\quad Count letters, test palindromes, or match a short pattern.\par
\entrylabel{Cost}\quad O(n) scans; O(nm) direct matching.\par
\entrylabel{Watch for}\quad Java char values are UTF-16 code units; the frequency helper accepts a-z only.\par
\end{guidebox}
\noindent\textbf{Core idea.} A string is an indexed sequence. Start with direct comparisons before choosing a faster string algorithm.\par
\noindent\textbf{Worked example.} The pattern aa occurs three times in aaaa because matches can overlap.\par
\noindent\textbf{Practice.} Test an empty text, one character, and overlapping matches.\par
\begin{lstlisting}
package notebook;

/** Character and substring patterns for introductory contest problems. */
public final class BeginnerStrings {
    private BeginnerStrings() {}

    /** Compare characters from both ends. O(n) time, O(1) extra space. */
    public static boolean isPalindrome(String text) {
        for (int left = 0, right = text.length() - 1; left < right; left++, right--)
            if (text.charAt(left) != text.charAt(right)) return false;
        return true;
    }

    /** Count lowercase English letters. O(n) time, O(1) space; rejects other characters. */
    public static int[] lowercaseFrequencies(String text) {
        int[] frequency = new int[26];
        for (int i = 0; i < text.length(); i++) {
            char ch = text.charAt(i);
            if (ch < 'a' || ch > 'z') throw new IllegalArgumentException("expected a-z");
            frequency[ch - 'a']++;
        }
        return frequency;
    }

    /** Count exact occurrences, including overlaps. O(n*m) time. */
    public static int countOccurrences(String text, String pattern) {
        if (pattern.isEmpty()) throw new IllegalArgumentException("pattern must not be empty");
        int count = 0;
        for (int start = 0; start + pattern.length() <= text.length(); start++) {
            int matched = 0;
            while (matched < pattern.length() && text.charAt(start + matched) == pattern.charAt(matched))
                matched++;
            if (matched == pattern.length()) count++;
        }
        return count;
    }

    /** Longest shared prefix of two strings. O(min(n,m)) time. */
    public static String commonPrefix(String first, String second) {
        int length = 0;
        while (length < first.length() && length < second.length()
                && first.charAt(length) == second.charAt(length)) length++;
        return first.substring(0, length);
    }
}
\end{lstlisting}
\Needspace{12\baselineskip}\subsection{KMP, Manacher, Boyer-Moore, Z}
\begin{guidebox}\small
\entrylabel{Level}\quad Intermediate\par
\entrylabel{Use when}\quad Search one pattern or find the longest palindrome.\par
\entrylabel{Cost}\quad Linear for KMP, Manacher, and Z.\par
\entrylabel{Watch for}\quad Uses Java UTF-16 code units.\par
\end{guidebox}
\noindent\textbf{Core idea.} Prefix, Z, and palindrome arrays summarize repeated character comparisons.\par
\noindent\textbf{Worked example.} The prefix function at i gives the longest proper border of the prefix ending at i.\par
\noindent\textbf{Practice.} Compare matches with direct substring search on random short strings.\par
\begin{lstlisting}
package notebook;

import java.util.ArrayList;
import java.util.Arrays;
import java.util.List;

/** String algorithms operate on Java UTF-16 code units. */
public final class StringsAlgorithms {
    private StringsAlgorithms() {}

    /** KMP match starts, including all n+1 positions for an empty pattern. O(n+m). */
    public static List<Integer> kmp(String text, String pattern) {
        List<Integer> matches = new ArrayList<>();
        if (pattern.isEmpty()) {
            for (int i = 0; i <= text.length(); i++) matches.add(i);
            return matches;
        }
        int[] pi = prefixFunction(pattern);
        for (int i = 0, j = 0; i < text.length(); i++) {
            while (j > 0 && text.charAt(i) != pattern.charAt(j)) j = pi[j - 1];
            if (text.charAt(i) == pattern.charAt(j)) j++;
            if (j == pattern.length()) { matches.add(i - j + 1); j = pi[j - 1]; }
        }
        return matches;
    }

    /** Prefix function. O(n). */
    public static int[] prefixFunction(String s) {
        int[] pi = new int[s.length()];
        for (int i = 1; i < s.length(); i++) {
            int j = pi[i - 1];
            while (j > 0 && s.charAt(i) != s.charAt(j)) j = pi[j - 1];
            if (s.charAt(i) == s.charAt(j)) j++;
            pi[i] = j;
        }
        return pi;
    }

    /** Longest palindromic substring by Manacher, O(n). */
    public static String manacher(String s) {
        if (s.isEmpty()) return "";
        int n = s.length(), bestStart = 0, bestLen = 1;
        int[] odd = new int[n], even = new int[n];
        for (int i = 0, l = 0, r = -1; i < n; i++) {
            int k = i > r ? 1 : Math.min(odd[l + r - i], r - i + 1);
            while (i - k >= 0 && i + k < n && s.charAt(i - k) == s.charAt(i + k)) k++;
            odd[i] = k--;
            int len = 2 * odd[i] - 1;
            if (len > bestLen) { bestLen = len; bestStart = i - odd[i] + 1; }
            if (i + k > r) { l = i - k; r = i + k; }
        }
        for (int i = 0, l = 0, r = -1; i < n; i++) {
            int k = i > r ? 0 : Math.min(even[l + r - i + 1], r - i + 1);
            while (i - k - 1 >= 0 && i + k < n && s.charAt(i - k - 1) == s.charAt(i + k)) k++;
            even[i] = k--;
            int len = 2 * even[i];
            if (len > bestLen) { bestLen = len; bestStart = i - even[i]; }
            if (i + k > r) { l = i - k - 1; r = i + k; }
        }
        return s.substring(bestStart, bestStart + bestLen);
    }

    /** Boyer-Moore bad-character search. O(nm) worst case; UTF-16 aware. */
    public static List<Integer> boyerMoore(String text, String pattern) {
        List<Integer> matches = new ArrayList<>();
        int n = text.length(), m = pattern.length();
        if (m == 0) {
            for (int i = 0; i <= n; i++) matches.add(i);
            return matches;
        }
        int[] last = new int[Character.MAX_VALUE + 1];
        Arrays.fill(last, -1);
        for (int i = 0; i < m; i++) last[pattern.charAt(i)] = i;
        for (int shift = 0; shift <= n - m;) {
            int j = m - 1;
            while (j >= 0 && pattern.charAt(j) == text.charAt(shift + j)) j--;
            if (j < 0) {
                matches.add(shift);
                shift += shift + m < n ? m - last[text.charAt(shift + m)] : 1;
            } else shift += Math.max(1, j - last[text.charAt(shift + j)]);
        }
        return matches;
    }

    /** Z-function: z[i] is longest common prefix of s and s.substring(i). O(n). */
    public static int[] zFunction(String s) {
        int[] z = new int[s.length()];
        for (int i = 1, l = 0, r = 0; i < s.length(); i++) {
            if (i <= r) z[i] = Math.min(r - i + 1, z[i - l]);
            while (i + z[i] < s.length() && s.charAt(z[i]) == s.charAt(i + z[i])) z[i]++;
            if (i + z[i] - 1 > r) { l = i; r = i + z[i] - 1; }
        }
        return z;
    }
}
\end{lstlisting}
\Needspace{12\baselineskip}\subsection{Trie and rolling hash}
\begin{guidebox}\small
\entrylabel{Level}\quad Intermediate\par
\entrylabel{Use when}\quad Store prefixes or compare substrings quickly.\par
\entrylabel{Cost}\quad O(total characters) build; O(1) hash query.\par
\entrylabel{Watch for}\quad Hash equality is probabilistic.\par
\end{guidebox}
\noindent\textbf{Core idea.} A trie shares common prefixes; rolling hash compresses a substring into modular values.\par
\noindent\textbf{Worked example.} Words car and cat share the c-a prefix nodes in a trie.\par
\noindent\textbf{Practice.} Use direct string comparison when a hash collision would be unacceptable.\par
\begin{lstlisting}
package notebook;

import java.util.HashMap;
import java.util.Map;

/** Trie and double rolling hash over Java UTF-16 code units. */
public final class AdvancedStrings {
    private AdvancedStrings() {}

    public static final class Trie {
        private static final class Node {
            Map<Character, Node> children = new HashMap<>();
            int terminalCount;
        }
        private final Node root = new Node();

        /** Insert a word, including the empty word. O(length). */
        public void add(String word) {
            Node node = root;
            for (int i = 0; i < word.length(); i++)
                node = node.children.computeIfAbsent(word.charAt(i), key -> new Node());
            node.terminalCount++;
        }

        public int count(String word) {
            Node node = root;
            for (int i = 0; i < word.length(); i++) {
                node = node.children.get(word.charAt(i));
                if (node == null) return 0;
            }
            return node.terminalCount;
        }

        public boolean hasPrefix(String prefix) {
            Node node = root;
            for (int i = 0; i < prefix.length(); i++) {
                node = node.children.get(prefix.charAt(i));
                if (node == null) return false;
            }
            return true;
        }
    }

    /** O(n) build, O(1) half-open substring hash. Hash equality is probabilistic. */
    public static final class RollingHash {
        private static final long MOD1 = 1_000_000_007L, MOD2 = 1_000_000_009L;
        private static final long BASE = 911_382_323L;
        private final long[] h1, h2, p1, p2;
        public RollingHash(String s) {
            int n = s.length();
            h1 = new long[n + 1]; h2 = new long[n + 1];
            p1 = new long[n + 1]; p2 = new long[n + 1];
            p1[0] = p2[0] = 1;
            for (int i = 0; i < n; i++) {
                long c = s.charAt(i) + 1L;
                h1[i + 1] = (h1[i] * BASE + c) % MOD1;
                h2[i + 1] = (h2[i] * BASE + c) % MOD2;
                p1[i + 1] = p1[i] * BASE % MOD1;
                p2[i + 1] = p2[i] * BASE % MOD2;
            }
        }
        public long hash(int left, int rightExclusive) {
            if (left < 0 || left > rightExclusive || rightExclusive >= h1.length)
                throw new IndexOutOfBoundsException();
            int len = rightExclusive - left;
            long a = (h1[rightExclusive] - h1[left] * p1[len] % MOD1 + MOD1) % MOD1;
            long b = (h2[rightExclusive] - h2[left] * p2[len] % MOD2 + MOD2) % MOD2;
            return (a << 32) | b;
        }
    }
}
\end{lstlisting}
\Needspace{12\baselineskip}\subsection{Minimum cyclic rotation}
\begin{guidebox}\small
\entrylabel{Level}\quad Intermediate\par
\entrylabel{Use when}\quad Find a lexicographically smallest rotation.\par
\entrylabel{Cost}\quad O(n) time.\par
\entrylabel{Watch for}\quad Uses Java UTF-16 code units.\par
\end{guidebox}
\noindent\textbf{Core idea.} Booth discards candidate starts after each mismatch.\par
\noindent\textbf{Worked example.} The minimum rotation of bba is abb.\par
\noindent\textbf{Practice.} Compare all rotations on short strings.\par
\begin{lstlisting}
package notebook;

/** Lexicographically minimum cyclic rotation by Booth's linear-time algorithm. */
public final class StringExtensions {
    private StringExtensions() {}

    /** Index of the smallest rotation; zero for empty text. O(n) time, O(1) extra space. */
    public static int minimumRotationIndex(String text) {
        int n = text.length();
        if (n == 0) return 0;
        int first = 0, second = 1, offset = 0;
        while (first < n && second < n && offset < n) {
            char a = text.charAt((first + offset) % n);
            char b = text.charAt((second + offset) % n);
            if (a == b) { offset++; continue; }
            if (a > b) first += offset + 1;
            else second += offset + 1;
            if (first == second) second++;
            offset = 0;
        }
        return Math.min(first, second);
    }

    /** Lexicographically smallest cyclic rotation. O(n). */
    public static String minimumRotation(String text) {
        int index = minimumRotationIndex(text);
        return text.substring(index) + text.substring(0, index);
    }
}
\end{lstlisting}
\Needspace{10\baselineskip}\section{Number theory and combinatorics}
\begin{guidebox}Use exact or modular arithmetic according to the requested result.\end{guidebox}
\Needspace{12\baselineskip}\subsection{GCD, trial division, divisors, and powers}
\begin{guidebox}\small
\entrylabel{Level}\quad Beginner\par
\entrylabel{Use when}\quad Solve small divisibility and integer arithmetic tasks.\par
\entrylabel{Cost}\quad GCD O(log n); trial division O(sqrt n); power O(log exponent).\par
\entrylabel{Watch for}\quad Exact power reports long overflow; inputs to gcd must be nonnegative.\par
\end{guidebox}
\noindent\textbf{Core idea.} A divisor larger than the square root has a matching smaller divisor. Euclid repeatedly replaces (a,b) by (b,a mod b).\par
\noindent\textbf{Worked example.} The divisors of 36 come in pairs around 6, giving nine positive divisors.\par
\noindent\textbf{Practice.} Compare trial-division results with direct enumeration for small values.\par
\begin{lstlisting}
package notebook;

/** Elementary integer math before modular arithmetic and number theory. */
public final class BeginnerMath {
    private BeginnerMath() {}

    /** Euclid's algorithm for nonnegative inputs. O(log min(a,b)) time. */
    public static long gcd(long a, long b) {
        if (a < 0 || b < 0) throw new IllegalArgumentException("expected nonnegative values");
        while (b != 0) { long next = a % b; a = b; b = next; }
        return a;
    }

    /** Test primality by trial division. O(sqrt(n)) time, O(1) space. */
    public static boolean isPrime(long value) {
        if (value < 2) return false;
        for (long divisor = 2; divisor <= value / divisor; divisor++)
            if (value % divisor == 0) return false;
        return true;
    }

    /** Count positive divisors of a positive integer. O(sqrt(n)) time. */
    public static long divisorCount(long value) {
        if (value <= 0) throw new IllegalArgumentException("expected a positive value");
        long count = 0;
        for (long divisor = 1; divisor <= value / divisor; divisor++)
            if (value % divisor == 0) count += divisor == value / divisor ? 1 : 2;
        return count;
    }

    /** Compute base^exponent for a nonnegative exponent. O(log exponent) time. */
    public static long power(long base, long exponent) {
        if (exponent < 0) throw new IllegalArgumentException("negative exponent");
        long result = 1;
        while (exponent != 0) {
            if ((exponent & 1) != 0) result = Math.multiplyExact(result, base);
            exponent >>= 1;
            if (exponent != 0) base = Math.multiplyExact(base, base);
        }
        return result;
    }
}
\end{lstlisting}
\Needspace{12\baselineskip}\subsection{GCD, modular arithmetic, primality, CRT}
\begin{guidebox}\small
\entrylabel{Level}\quad Intermediate\par
\entrylabel{Use when}\quad Work with divisibility, primes, congruences, or roots.\par
\entrylabel{Cost}\quad Euclid O(log n); sieve O(n log log n).\par
\entrylabel{Watch for}\quad Overflow-safe modular multiplication trades speed for range.\par
\end{guidebox}
\noindent\textbf{Core idea.} Choose the correct number representation and modulus before applying arithmetic routines.\par
\noindent\textbf{Worked example.} Euclid finds gcd; the extended form also gives coefficients for modular inverses.\par
\noindent\textbf{Practice.} Compare modular operations against BigInteger on boundary inputs.\par
\begin{lstlisting}
package notebook;

import java.math.BigInteger;
import java.util.ArrayList;
import java.util.Collections;
import java.util.List;
import java.util.function.DoubleUnaryOperator;

/** Integer number theory. Modular operations require mod > 0. */
public final class NumberTheory {
    private NumberTheory() {}

    public static long gcd(long a, long b) {
        if (a == Long.MIN_VALUE || b == Long.MIN_VALUE)
            return BigInteger.valueOf(a).gcd(BigInteger.valueOf(b)).longValueExact();
        a = Math.abs(a); b = Math.abs(b);
        while (b != 0) { long t = a % b; a = b; b = t; }
        return a;
    }

    public static long lcm(long a, long b) {
        if (a == 0 || b == 0) return 0;
        long product = Math.multiplyExact(a / gcd(a, b), b);
        if (product == Long.MIN_VALUE) throw new ArithmeticException("lcm does not fit in long");
        return Math.abs(product);
    }

    /** Returns [gcd(a,b), x, y] with ax+by=gcd; coefficients must fit in long. */
    public static long[] extendedGcd(long a, long b) {
        BigInteger aa = BigInteger.valueOf(a), bb = BigInteger.valueOf(b);
        BigInteger oldR = aa, r = bb, oldS = BigInteger.ONE, s = BigInteger.ZERO;
        BigInteger oldT = BigInteger.ZERO, t = BigInteger.ONE;
        while (r.signum() != 0) {
            BigInteger q = oldR.divide(r);
            BigInteger nr = oldR.subtract(q.multiply(r)); oldR = r; r = nr;
            BigInteger ns = oldS.subtract(q.multiply(s)); oldS = s; s = ns;
            BigInteger nt = oldT.subtract(q.multiply(t)); oldT = t; t = nt;
        }
        if (oldR.signum() < 0) { oldR = oldR.negate(); oldS = oldS.negate(); oldT = oldT.negate(); }
        return new long[] { oldR.longValueExact(), oldS.longValueExact(), oldT.longValueExact() };
    }

    /** Exact multiplication modulo mod without signed-long overflow. O(log b). */
    public static long modMul(long a, long b, long mod) {
        if (mod <= 0) throw new IllegalArgumentException("mod must be positive");
        long x = Math.floorMod(a, mod), y = Math.floorMod(b, mod), result = 0;
        while (y > 0) {
            if ((y & 1) != 0) result = addMod(result, x, mod);
            y >>>= 1;
            if (y > 0) x = addMod(x, x, mod);
        }
        return result;
    }

    private static long addMod(long a, long b, long mod) {
        return a >= mod - b ? a - (mod - b) : a + b;
    }

    /** a^exponent modulo mod. O(log exponent * log mod). */
    public static long modPow(long a, long exponent, long mod) {
        if (exponent < 0 || mod <= 0) throw new IllegalArgumentException();
        long result = 1 % mod, base = Math.floorMod(a, mod);
        while (exponent > 0) {
            if ((exponent & 1) != 0) result = modMul(result, base, mod);
            exponent >>>= 1;
            if (exponent > 0) base = modMul(base, base, mod);
        }
        return result;
    }

    /** Deterministic Miller-Rabin for signed 64-bit positive integers. */
    public static boolean isPrime(long n) {
        if (n < 2) return false;
        for (long p : new long[] {2, 3, 5, 7, 11, 13, 17, 19, 23, 29, 31, 37}) {
            if (n % p == 0) return n == p;
        }
        long d = n - 1;
        int s = Long.numberOfTrailingZeros(d);
        d >>= s;
        for (long witness : new long[] {2, 325, 9375, 28178, 450775, 9780504, 1795265022}) {
            long a = witness % n;
            if (a == 0) continue;
            long x = modPow(a, d, n);
            if (x == 1 || x == n - 1) continue;
            boolean passed = false;
            for (int r = 1; r < s; r++) {
                x = modMul(x, x, n);
                if (x == n - 1) { passed = true; break; }
            }
            if (!passed) return false;
        }
        return true;
    }

    /** Primes <= n; O(n log log n) time, O(n) space. */
    public static List<Integer> sieve(int n) {
        if (n < 2) return new ArrayList<>();
        boolean[] composite = new boolean[n + 1];
        for (int p = 2; (long) p * p <= n; p++)
            if (!composite[p]) for (int x = p * p; x <= n; x += p) composite[x] = true;
        List<Integer> primes = new ArrayList<>();
        for (int i = 2; i <= n; i++) if (!composite[i]) primes.add(i);
        return primes;
    }

    /** Sorted positive divisors of n; n > 0. O(sqrt n). */
    public static List<Long> divisors(long n) {
        if (n <= 0) throw new IllegalArgumentException();
        List<Long> result = new ArrayList<>();
        for (long d = 1; d <= n / d; d++) {
            if (n % d == 0) {
                result.add(d);
                if (d != n / d) result.add(n / d);
            }
        }
        Collections.sort(result);
        return result;
    }

    /** Exact factorial; n >= 0. */
    public static BigInteger factorial(int n) {
        if (n < 0) throw new IllegalArgumentException();
        BigInteger result = BigInteger.ONE;
        for (int i = 2; i <= n; i++) result = result.multiply(BigInteger.valueOf(i));
        return result;
    }

    /** Generalized CRT. Returns [least nonnegative solution, lcm], or null if inconsistent. */
    public static BigInteger[] chineseRemainder(long[] moduli, long[] residues) {
        if (moduli.length != residues.length) throw new IllegalArgumentException();
        BigInteger x = BigInteger.ZERO, m = BigInteger.ONE;
        for (int i = 0; i < moduli.length; i++) {
            if (moduli[i] <= 0) throw new IllegalArgumentException("moduli must be positive");
            BigInteger next = BigInteger.valueOf(moduli[i]);
            BigInteger target = BigInteger.valueOf(residues[i]).mod(next);
            BigInteger g = m.gcd(next), delta = target.subtract(x);
            if (!delta.mod(g).equals(BigInteger.ZERO)) return null;
            BigInteger reduced = next.divide(g);
            BigInteger k = reduced.equals(BigInteger.ONE) ? BigInteger.ZERO
                : delta.divide(g).multiply(m.divide(g).modInverse(reduced)).mod(reduced);
            x = x.add(m.multiply(k));
            m = m.multiply(reduced);
            x = x.mod(m);
        }
        return new BigInteger[] { x, m };
    }

    /** Newton iteration; derivative must remain nonzero near the root. */
    public static double newtonRaphson(DoubleUnaryOperator f, DoubleUnaryOperator derivative,
                                        double start, double tolerance, int maxIterations) {
        if (!(tolerance > 0) || maxIterations <= 0) throw new IllegalArgumentException();
        double x = start;
        for (int i = 0; i < maxIterations; i++) {
            double fx = f.applyAsDouble(x), dx = derivative.applyAsDouble(x);
            if (Math.abs(fx) <= tolerance) return x;
            if (!Double.isFinite(dx) || dx == 0) throw new ArithmeticException("zero derivative");
            x -= fx / dx;
            if (!Double.isFinite(x)) throw new ArithmeticException("non-finite iterate");
        }
        throw new ArithmeticException("did not converge");
    }
}
\end{lstlisting}
\Needspace{12\baselineskip}\subsection{Exact binomial and Catalan numbers}
\begin{guidebox}\small
\entrylabel{Level}\quad Intermediate\par
\entrylabel{Use when}\quad Compute counts beyond long range.\par
\entrylabel{Cost}\quad O(n) BigInteger operations.\par
\entrylabel{Watch for}\quad BigInteger is exact but allocates memory.\par
\end{guidebox}
\noindent\textbf{Core idea.} BigInteger keeps exact counts when the answer does not fit a primitive type.\par
\noindent\textbf{Worked example.} The number of ways to choose 2 from 5 is 10.\par
\noindent\textbf{Practice.} Check symmetry C(n,k) = C(n,n-k).\par
\begin{lstlisting}
package notebook;

import java.math.BigInteger;

public final class Combinatorics {
    private Combinatorics() {}

    /** Exact binomial coefficient. O(min(r,n-r)) big-integer operations. */
    public static BigInteger binomial(int n, int r) {
        if (r < 0 || r > n) return BigInteger.ZERO;
        r = Math.min(r, n - r);
        BigInteger ans = BigInteger.ONE;
        for (int i = 1; i <= r; i++)
            ans = ans.multiply(BigInteger.valueOf(n - i + 1)).divide(BigInteger.valueOf(i));
        return ans;
    }

    /** Catalan numbers C[0..n], exact. O(n) big-integer operations. */
    public static BigInteger[] catalan(int n) {
        if (n < 0) throw new IllegalArgumentException();
        BigInteger[] c = new BigInteger[n + 1];
        c[0] = BigInteger.ONE;
        for (int i = 1; i <= n; i++)
            c[i] = c[i - 1].multiply(BigInteger.valueOf(4L * i - 2)).divide(BigInteger.valueOf(i + 1));
        return c;
    }
}
\end{lstlisting}
\Needspace{12\baselineskip}\subsection{Derangements, Stirling, partitions, Burnside}
\begin{guidebox}\small
\entrylabel{Level}\quad Intermediate\par
\entrylabel{Use when}\quad Count permutations, set partitions, and necklaces.\par
\entrylabel{Cost}\quad O(n) to O(n sqrt n) BigInteger work.\par
\entrylabel{Watch for}\quad Exact answers can become very large.\par
\end{guidebox}
\noindent\textbf{Core idea.} Recurrences and Burnside count objects without enumerating them.\par
\noindent\textbf{Worked example.} There are 44 derangements of five items.\par
\noindent\textbf{Practice.} Enumerate tiny objects and compare counts.\par
\begin{lstlisting}
package notebook;

import java.math.BigInteger;

/** Exact counting formulas for permutations, partitions, and rotational symmetry. */
public final class CombinatoricsExtras {
    private CombinatoricsExtras() {}

    /** Derangements of n labeled items. O(n) BigInteger operations. */
    public static BigInteger derangements(int n) {
        if (n < 0) throw new IllegalArgumentException();
        if (n == 0) return BigInteger.ONE;
        BigInteger previous = BigInteger.ONE, current = BigInteger.ZERO;
        for (int i = 2; i <= n; i++) {
            BigInteger next = previous.add(current).multiply(BigInteger.valueOf(i - 1));
            previous = current; current = next;
        }
        return current;
    }

    /** Stirling number S(n,k): partitions into exactly k nonempty unlabeled groups. O(n*k). */
    public static BigInteger stirlingSecond(int n, int k) {
        if (n < 0 || k < 0) throw new IllegalArgumentException();
        if (k > n) return BigInteger.ZERO;
        BigInteger[] dp = new BigInteger[k + 1];
        for (int i = 0; i <= k; i++) dp[i] = BigInteger.ZERO;
        dp[0] = BigInteger.ONE;
        for (int value = 1; value <= n; value++)
            for (int groups = Math.min(value, k); groups >= 0; groups--)
                dp[groups] = groups == 0 ? BigInteger.ZERO
                        : dp[groups - 1].add(dp[groups].multiply(BigInteger.valueOf(groups)));
        return dp[k];
    }

    /** Partition numbers p(0)..p(limit) using generalized pentagonal numbers. O(n*sqrt(n)). */
    public static BigInteger[] partitions(int limit) {
        if (limit < 0) throw new IllegalArgumentException();
        BigInteger[] result = new BigInteger[limit + 1];
        result[0] = BigInteger.ONE;
        for (int n = 1; n <= limit; n++) {
            BigInteger sum = BigInteger.ZERO;
            for (long k = 1; k * (3 * k - 1) / 2 <= n; k++) {
                int first = (int) (k * (3 * k - 1) / 2);
                int second = (int) (k * (3 * k + 1) / 2);
                if ((k & 1) == 1) {
                    sum = sum.add(result[n - first]);
                    if (second <= n) sum = sum.add(result[n - second]);
                } else {
                    sum = sum.subtract(result[n - first]);
                    if (second <= n) sum = sum.subtract(result[n - second]);
                }
            }
            result[n] = sum;
        }
        return result;
    }

    /** Number of length-n necklaces over colors colors, up to rotation. Burnside in O(n log n). */
    public static BigInteger necklaces(int n, int colors) {
        if (n <= 0 || colors < 0) throw new IllegalArgumentException();
        BigInteger sum = BigInteger.ZERO, base = BigInteger.valueOf(colors);
        for (int shift = 0; shift < n; shift++) sum = sum.add(base.pow(gcd(n, shift)));
        return sum.divide(BigInteger.valueOf(n));
    }

    private static int gcd(int a, int b) {
        while (b != 0) { int next = a % b; a = b; b = next; }
        return a;
    }
}
\end{lstlisting}
\Needspace{10\baselineskip}\section{Geometry}
\begin{guidebox}Use floating-point distances and exact integer predicates as appropriate.\end{guidebox}
\Needspace{12\baselineskip}\subsection{Distance, polygons, hull, closest pair}
\begin{guidebox}\small
\entrylabel{Level}\quad Intermediate\par
\entrylabel{Use when}\quad Measure and optimize among planar points.\par
\entrylabel{Cost}\quad O(n) to O(n log\textasciicircum{}2 n).\par
\entrylabel{Watch for}\quad Floating-point orientation is unstable near collinearity.\par
\end{guidebox}
\noindent\textbf{Core idea.} Sort and orient points carefully; choose exact predicates when coordinates are integral.\par
\noindent\textbf{Worked example.} A convex hull removes interior points before farthest-pair search.\par
\noindent\textbf{Practice.} Test duplicate and collinear points.\par
\begin{lstlisting}
package notebook;

import java.util.ArrayList;
import java.util.Arrays;
import java.util.Comparator;
import java.util.List;

/** Planar geometry with double coordinates; use exact arithmetic for integer-critical predicates. */
public final class Geometry {
    private Geometry() {}

    public record Point(double x, double y) {}

    public static double distance(Point a, Point b) {
        return Math.hypot(a.x - b.x, a.y - b.y);
    }

    /** Shoelace area; input vertices are not repeated at the end. O(n). */
    public static double polygonArea(List<Point> p) {
        double twice = 0;
        for (int i = 0; i < p.size(); i++) {
            Point a = p.get(i), b = p.get((i + 1) % p.size());
            twice += a.x * b.y - a.y * b.x;
        }
        return Math.abs(twice) / 2;
    }

    /** Closed polygon perimeter; O(n). */
    public static double polygonPerimeter(List<Point> p) {
        double result = 0;
        for (int i = 0; i < p.size(); i++) result += distance(p.get(i), p.get((i + 1) % p.size()));
        return result;
    }

    private static double cross(Point a, Point b, Point c) {
        return (b.x - a.x) * (c.y - a.y) - (b.y - a.y) * (c.x - a.x);
    }

    /** Monotone-chain hull in counterclockwise order; removes collinear interior points. O(n log n). */
    public static List<Point> convexHull(Point[] points) {
        Point[] sorted = points.clone();
        Arrays.sort(sorted, Comparator.comparingDouble(Point::x).thenComparingDouble(Point::y));
        List<Point> unique = new ArrayList<>();
        for (Point p : sorted) if (unique.isEmpty() || !unique.get(unique.size() - 1).equals(p)) unique.add(p);
        if (unique.size() <= 1) return unique;
        List<Point> lower = new ArrayList<>(), upper = new ArrayList<>();
        for (Point p : unique) {
            while (lower.size() >= 2 && cross(lower.get(lower.size() - 2), lower.get(lower.size() - 1), p) <= 0)
                lower.remove(lower.size() - 1);
            lower.add(p);
        }
        for (int i = unique.size() - 1; i >= 0; i--) {
            Point p = unique.get(i);
            while (upper.size() >= 2 && cross(upper.get(upper.size() - 2), upper.get(upper.size() - 1), p) <= 0)
                upper.remove(upper.size() - 1);
            upper.add(p);
        }
        lower.remove(lower.size() - 1);
        upper.remove(upper.size() - 1);
        lower.addAll(upper);
        return lower;
    }

    /** Closest pair distance. O(n log^2 n); requires at least two points. */
    public static double closestPair(Point[] points) {
        if (points.length < 2) throw new IllegalArgumentException("at least two points required");
        Point[] sorted = points.clone();
        Arrays.sort(sorted, Comparator.comparingDouble(Point::x).thenComparingDouble(Point::y));
        return closest(sorted, 0, sorted.length);
    }

    private static double closest(Point[] p, int lo, int hi) {
        if (hi - lo <= 3) {
            double best = Double.POSITIVE_INFINITY;
            for (int i = lo; i < hi; i++) for (int j = i + 1; j < hi; j++) best = Math.min(best, distance(p[i], p[j]));
            return best;
        }
        int mid = (lo + hi) >>> 1;
        double splitX = p[mid].x;
        double best = Math.min(closest(p, lo, mid), closest(p, mid, hi));
        List<Point> strip = new ArrayList<>();
        for (int i = lo; i < hi; i++) if (Math.abs(p[i].x - splitX) < best) strip.add(p[i]);
        strip.sort(Comparator.comparingDouble(Point::y));
        for (int i = 0; i < strip.size(); i++)
            for (int j = i + 1; j < strip.size() && strip.get(j).y - strip.get(i).y < best; j++)
                best = Math.min(best, distance(strip.get(i), strip.get(j)));
        return best;
    }

    /** Diameter via convex hull and rotating calipers. O(n log n). */
    public static double farthestPair(Point[] points) {
        List<Point> hull = convexHull(points);
        int n = hull.size();
        if (n < 2) return 0;
        if (n == 2) return distance(hull.get(0), hull.get(1));
        int j = 1;
        double best = 0;
        for (int i = 0; i < n; i++) {
            int next = (i + 1) % n;
            while (Math.abs(cross(hull.get(i), hull.get(next), hull.get((j + 1) % n)))
                    > Math.abs(cross(hull.get(i), hull.get(next), hull.get(j)))) j = (j + 1) % n;
            best = Math.max(best, Math.max(distance(hull.get(i), hull.get(j)), distance(hull.get(next), hull.get(j))));
        }
        return best;
    }
}
\end{lstlisting}
\Needspace{12\baselineskip}\subsection{Lines, segments, and circle intersections}
\begin{guidebox}\small
\entrylabel{Level}\quad Intermediate\par
\entrylabel{Use when}\quad Measure distances or intersect lines and circles.\par
\entrylabel{Cost}\quad O(1) per primitive.\par
\entrylabel{Watch for}\quad Coincident circles have infinitely many intersections.\par
\end{guidebox}
\noindent\textbf{Core idea.} Projection and the common chord solve these primitives.\par
\noindent\textbf{Worked example.} Unit circles centered at (0,0) and (2,0) touch at (1,0).\par
\noindent\textbf{Practice.} Test tangent, disjoint, and degenerate cases.\par
\begin{lstlisting}
package notebook;

import java.util.List;

/** Floating-point line, segment, and circle primitives with an absolute epsilon. */
public final class CircleGeometry {
    private CircleGeometry() {}
    private static final double EPS = 1e-10;
    public record Point(double x, double y) {}

    /** Intersection of two infinite lines, or null when parallel/coincident. O(1). */
    public static Point lineIntersection(Point a, Point b, Point c, Point d) {
        double ux = b.x - a.x, uy = b.y - a.y;
        double vx = d.x - c.x, vy = d.y - c.y;
        double determinant = ux * vy - uy * vx;
        if (Math.abs(determinant) < EPS) return null;
        double t = ((c.x - a.x) * vy - (c.y - a.y) * vx) / determinant;
        return new Point(a.x + t * ux, a.y + t * uy);
    }

    /** Euclidean distance from a point to a closed segment. O(1). */
    public static double pointSegmentDistance(Point point, Point a, Point b) {
        double dx = b.x - a.x, dy = b.y - a.y;
        double squared = dx * dx + dy * dy;
        if (squared == 0) return Math.hypot(point.x - a.x, point.y - a.y);
        double t = Math.max(0, Math.min(1, ((point.x - a.x) * dx + (point.y - a.y) * dy) / squared));
        return Math.hypot(point.x - (a.x + t * dx), point.y - (a.y + t * dy));
    }

    /** Intersections of an infinite line with a circle. Returns zero, one, or two points. */
    public static List<Point> circleLine(Point center, double radius, Point a, Point b) {
        if (radius < 0) throw new IllegalArgumentException("negative radius");
        double dx = b.x - a.x, dy = b.y - a.y;
        double lengthSquared = dx * dx + dy * dy;
        if (lengthSquared == 0) throw new IllegalArgumentException("line needs two distinct points");
        double t = ((center.x - a.x) * dx + (center.y - a.y) * dy) / lengthSquared;
        Point foot = new Point(a.x + t * dx, a.y + t * dy);
        double offsetSquared = radius * radius - Math.pow(foot.x - center.x, 2) - Math.pow(foot.y - center.y, 2);
        if (offsetSquared < -EPS) return List.of();
        if (Math.abs(offsetSquared) <= EPS) return List.of(foot);
        double offset = Math.sqrt(offsetSquared / lengthSquared);
        return List.of(new Point(foot.x + offset * dx, foot.y + offset * dy),
                new Point(foot.x - offset * dx, foot.y - offset * dy));
    }

    /** Intersections of two circles; coincident circles are rejected. */
    public static List<Point> circleCircle(Point first, double firstRadius, Point second, double secondRadius) {
        if (firstRadius < 0 || secondRadius < 0) throw new IllegalArgumentException("negative radius");
        double dx = second.x - first.x, dy = second.y - first.y;
        double distance = Math.hypot(dx, dy);
        if (distance < EPS) {
            if (Math.abs(firstRadius - secondRadius) < EPS)
                throw new IllegalArgumentException("infinitely many intersections");
            return List.of();
        }
        if (distance > firstRadius + secondRadius + EPS
                || distance < Math.abs(firstRadius - secondRadius) - EPS) return List.of();
        double along = (firstRadius * firstRadius - secondRadius * secondRadius + distance * distance)
                / (2 * distance);
        double heightSquared = firstRadius * firstRadius - along * along;
        if (heightSquared < -EPS) return List.of();
        double x = first.x + along * dx / distance, y = first.y + along * dy / distance;
        if (Math.abs(heightSquared) <= EPS) return List.of(new Point(x, y));
        double height = Math.sqrt(heightSquared);
        return List.of(new Point(x - height * dy / distance, y + height * dx / distance),
                new Point(x + height * dy / distance, y - height * dx / distance));
    }
}
\end{lstlisting}
\Needspace{12\baselineskip}\subsection{3D vectors and spherical distance}
\begin{guidebox}\small
\entrylabel{Level}\quad Intermediate\par
\entrylabel{Use when}\quad Compute vector products, tetrahedron volume, or distance on a sphere.\par
\entrylabel{Cost}\quad O(1) per operation.\par
\entrylabel{Watch for}\quad Latitude uses degrees; spherical distance assumes a mean Earth radius.\par
\end{guidebox}
\noindent\textbf{Core idea.} Dot and cross products extend orientation and area calculations to three dimensions.\par
\noindent\textbf{Worked example.} The cross product of (1,0,0) and (0,1,0) is (0,0,1).\par
\noindent\textbf{Practice.} Check orthogonality with dot products and compare great-circle distances with known locations.\par
\begin{lstlisting}
package notebook;

/** Three-dimensional vector operations and spherical distance. */
public final class Geometry3D {
    private Geometry3D() {}
    public record Vector(double x, double y, double z) {
        public Vector add(Vector other) { return new Vector(x+other.x,y+other.y,z+other.z); }
        public Vector subtract(Vector other) { return new Vector(x-other.x,y-other.y,z-other.z); }
        public Vector scale(double factor) { return new Vector(x*factor,y*factor,z*factor); }
        public double dot(Vector other) { return x*other.x+y*other.y+z*other.z; }
        public Vector cross(Vector other) {
            return new Vector(y*other.z-z*other.y,z*other.x-x*other.z,x*other.y-y*other.x);
        }
        public double norm() { return Math.sqrt(dot(this)); }
        public double distance(Vector other) { return subtract(other).norm(); }
    }

    /** Signed volume of tetrahedron (origin,a,b,c). */
    public static double signedTetrahedronVolume(Vector a, Vector b, Vector c) {
        return a.dot(b.cross(c))/6.0;
    }

    /** Great-circle distance in kilometers using the haversine formula. */
    public static double greatCircleKilometers(double latitude1, double longitude1,
                                                double latitude2, double longitude2) {
        if (Math.abs(latitude1)>90 || Math.abs(latitude2)>90) throw new IllegalArgumentException("invalid latitude");
        double lat1=Math.toRadians(latitude1), lat2=Math.toRadians(latitude2);
        double dLat=lat2-lat1, dLon=Math.toRadians(longitude2-longitude1);
        double haversine=Math.sin(dLat/2)*Math.sin(dLat/2)
                +Math.cos(lat1)*Math.cos(lat2)*Math.sin(dLon/2)*Math.sin(dLon/2);
        return 6371.0088*2*Math.asin(Math.sqrt(Math.max(0,Math.min(1,haversine))));
    }
}
\end{lstlisting}
\Needspace{10\baselineskip}\section{Game theory}
\begin{guidebox}Model legal moves and terminal states before applying xor or mex.\end{guidebox}
\Needspace{12\baselineskip}\subsection{Nim, mex, Grundy, minimax}
\begin{guidebox}\small
\entrylabel{Level}\quad Intermediate\par
\entrylabel{Use when}\quad Analyze impartial games or complete binary game trees.\par
\entrylabel{Cost}\quad O(piles), O(edges), or O(leaves).\par
\entrylabel{Watch for}\quad Nim assumes normal play.\par
\end{guidebox}
\noindent\textbf{Core idea.} For impartial normal-play games, a zero xor of Grundy values is losing.\par
\noindent\textbf{Worked example.} Nim heaps of sizes 1, 2, and 3 have xor zero.\par
\noindent\textbf{Practice.} Enumerate all moves in a tiny game to verify Grundy labels.\par
\begin{lstlisting}
package notebook;

import java.util.HashSet;
import java.util.Set;

public final class GameTheory {
    private GameTheory() {}

    /** Normal-play Nim: first player wins iff pile xor is nonzero. */
    public static boolean firstWinsNim(int[] piles) {
        int xor = 0;
        for (int pile : piles) {
            if (pile < 0) throw new IllegalArgumentException();
            xor ^= pile;
        }
        return xor != 0;
    }

    /** Minimum excluded nonnegative integer. */
    public static int mex(Set<Integer> values) {
        int result = 0;
        while (values.contains(result)) result++;
        return result;
    }

    /** Grundy values for a DAG whose edges only lead to lower-numbered states. */
    public static int[] grundy(int[][] moves) {
        int[] result = new int[moves.length];
        for (int i = 0; i < moves.length; i++) {
            Set<Integer> reachable = new HashSet<>();
            for (int to : moves[i]) {
                if (to < 0 || to >= i) throw new IllegalArgumentException("moves must lead to lower states");
                reachable.add(result[to]);
            }
            result[i] = mex(reachable);
        }
        return result;
    }

    /** Minimax of a complete binary game tree; leaves length must be a power of two. */
    public static int minimax(int[] leaves) {
        if (leaves.length == 0 || (leaves.length & (leaves.length - 1)) != 0)
            throw new IllegalArgumentException("leaf count must be a power of two");
        int[] level = leaves.clone();
        boolean maximize = (Integer.numberOfTrailingZeros(leaves.length) & 1) == 1;
        while (level.length > 1) {
            int[] next = new int[level.length / 2];
            for (int i = 0; i < next.length; i++)
                next[i] = maximize ? Math.max(level[2 * i], level[2 * i + 1])
                    : Math.min(level[2 * i], level[2 * i + 1]);
            level = next;
            maximize = !maximize;
        }
        return level[0];
    }
}
\end{lstlisting}
\end{document}
